\documentclass[11pt]{article}

\usepackage[a4paper,margin=22mm,headheight=14pt]{geometry}
\usepackage{fontspec}
\usepackage{xeCJK}
\xeCJKsetup{CJKspace=true}
\setmainfont{DejaVu Serif}
\setsansfont{DejaVu Sans}
\setmonofont{DejaVu Sans Mono}
\setCJKmainfont{Noto Serif CJK KR}
\setCJKsansfont{Noto Sans CJK KR}
\setCJKmonofont{Noto Sans Mono CJK KR}

\usepackage{amsmath,amssymb,bm,mathtools}
\usepackage{booktabs,longtable,tabularx,array}
\usepackage{enumitem}
\usepackage{xcolor}
\usepackage{hyperref}
\usepackage{fancyhdr}
\usepackage{microtype}
\usepackage{tocloft}
\usepackage[most]{tcolorbox}

\definecolor{CoreBlue}{HTML}{1F4E79}
\definecolor{CoreLight}{HTML}{EAF3F8}
\definecolor{LocalPurple}{HTML}{674EA7}
\definecolor{LocalLight}{HTML}{EEEAF7}
\definecolor{RiskOrange}{HTML}{B45F06}
\definecolor{RiskLight}{HTML}{FFF2CC}
\definecolor{StopRed}{HTML}{9C0006}
\definecolor{StopLight}{HTML}{FCE8E6}
\definecolor{SafeGreen}{HTML}{38761D}
\definecolor{SafeLight}{HTML}{EAF4E2}

\hypersetup{
  colorlinks=true,
  linkcolor=CoreBlue,
  citecolor=CoreBlue,
  urlcolor=CoreBlue,
  pdftitle={Response-Distilled Global-Local Wind Dynamics for Static 3D Gaussian Thin Surfaces},
  pdfauthor={Wind3DGS Project}
}

\pagestyle{fancy}
\fancyhf{}
\lhead{\small Response-Distilled Global--Local Wind Dynamics}
\rhead{\small Canonical Research Sketch}
\cfoot{\thepage}

\setlength{\parindent}{0pt}
\setlength{\parskip}{0.45em}
\setlist[itemize]{leftmargin=1.5em,itemsep=0.15em,topsep=0.25em}
\setlist[enumerate]{leftmargin=1.8em,itemsep=0.2em,topsep=0.25em}
\setlength{\cftsecnumwidth}{3em}
\setlength{\cftsubsecnumwidth}{4em}
\renewcommand{\arraystretch}{1.18}
\newcolumntype{L}[1]{>{\raggedright\arraybackslash}p{#1}}
\newcommand{\code}[1]{{\small\nolinkurl{#1}}}
\newcommand{\globaltag}[1]{\textcolor{CoreBlue}{\textbf{#1}}}
\newcommand{\localtag}[1]{\textcolor{LocalPurple}{\textbf{#1}}}
\newcommand{\risktag}[1]{\textcolor{RiskOrange}{\textbf{#1}}}
\newcommand{\stoptag}[1]{\textcolor{StopRed}{\textbf{#1}}}
\newcommand{\safetag}[1]{\textcolor{SafeGreen}{\textbf{#1}}}

\title{\textbf{Response-Distilled Global--Local Wind Dynamics}\\
for Static 3D Gaussian Thin Surfaces\\[0.35em]
\large Target-Mesh-Free Inference with Topology-Privileged Training}
\author{Wind3DGS Project}
\date{2026-08-22}

\begin{document}
\maketitle

\begin{tcolorbox}[colback=CoreLight,colframe=CoreBlue,title=Canonical direction as of 2026-08-22,fonttitle=\bfseries]
이 문서는 Wind3DGS의 현재 연구 질문, 방법 경계, 학습 데이터 계약과 구현 Gate를 소유한다.
이전 Minimal V0의 explicit anchor scaffold, decoded edge/triplet topology, fixed-Hessian 구조 연산자와
generalized-eigen Global/patch-complement Local solver는 복구 가능한 archive와 비교 baseline으로 남기되,
새 방법의 완료 상태로 승계하지 않는다.

새 방법은 training-only mesh simulation에서 \textbf{force-to-motion response}를 학습하고,
target에서는 static 3DGS, metric scale, material, attachment와 prescribed wind만 사용한다.
네트워크 내부의 patch token과 latent relation은 사용하지만 target mesh나 persistent physical graph를
외부 산출물로 만들지 않는다. 네트워크의 method output은 다음 frame의 absolute deformation이 아니라
Global 및 Local \textbf{response package}이며, 작은 deterministic response evaluator가 이를 시간 적분한다.
목표는 universal physical world model이 아니라 지원한 thin-surface domain에서
시각적으로 설득력 있고 안정적이며 controllable한 wind animation과 measured quality--latency 이득이다.
\end{tcolorbox}

\tableofcontents
\newpage

\section{한눈에 보는 아이디어}

\subsection{핵심 흐름}

\begin{tcolorbox}[colback=SafeLight,colframe=SafeGreen,title=Training,fonttitle=\bfseries]
연결정보가 있는 mesh로 여러 재료, 부착과 바람 조건의 reference simulation을 만든다.
물체별로 공간·시간 수렴을 확인해 선택한 teacher 해상도를 사용하고, 같은 물리 조건의 응답을
여러 GS 표현이 공유한다. 같은 rest surface를 GS 생성 시 Gaussian 예산 또는 densification 제한을
달리하여 독립적으로 reconstruction한다. Seed와 view 변화는 별도의 재구성 일반화 축으로 다룬다.
Mesh vertex와 topology는 teacher quadrature, same-surface relation과 response label을 만드는 데만 사용한다.
Student는 static 3DGS만 보고 density가 달라도 비슷한 내부 latent token을 형성하고,
그 token에서 passive Global response와 spatially localized Local residual response를 예측하도록 학습한다.
\end{tcolorbox}

\begin{tcolorbox}[colback=CoreLight,colframe=CoreBlue,title=Target setup and runtime,fonttitle=\bfseries]
Target static 3DGS와 사용자가 준 material/attachment를 setup network에 한 번 넣어 response package를 cache한다.
매 frame에는 바람이 만드는 distributed load를 package에 투영하고, Global response는 항상 적분한다.
현재 load에서 필요한 Local slot만 fixed budget 안에서 활성화해 flutter와 국소 phase detail을 보완한다.
결과 state로 Gaussian mean과 covariance orientation을 움직여 즉시 렌더한다.
Target mesh, target dynamic video, per-target physical fitting과 persistent edge graph는 요구하지 않는다.
\end{tcolorbox}

\paragraph{사용자 상호작용과 실행 시간 목표.}
사용자는 지원 범위의 바람 방향·세기를 조작하며 최소30fps로 반응을 보는 것을 목표로 한다.
힘 계산, 작은 response state 갱신, Gaussian 변형, 렌더링과 입력/동기화를 포함해 약33.3 ms/frame의
예산 안에 들어가는지 측정해야 한다. 오프라인 Teacher 생성·학습과 물체별 setup 비용은 별도로 보고한다.
이는 아직 실측하지 않은 목표이며, 큰 Teacher \(\Delta t\)를 찾거나 GPU로 계산한 사실만으로 달성되지 않는다.
Setup-once evaluator는 현재 변형·속도를 이어 바람에 반응하지만, 임의 잡기·충돌·고정점 이동은 현재 계약 밖이다.
고정 response field가 큰 변형을 충분히 표현하는지와 렌더링까지 포함한 비용은 각각 oracle/학습·profiler에서 검증한다.

\subsection{네 가지 표현을 혼동하지 않는다}

\begin{longtable}{L{0.20\linewidth}L{0.70\linewidth}}
\toprule
표현 & 역할\\
\midrule
\endfirsthead
\toprule
표현 & 역할\\
\midrule
\endhead
Rendering Gaussian & 입력 및 최종 렌더 primitive. 수와 density는 reconstruction마다 달라질 수 있다.\\
Latent response token & Patch evidence와 global context를 모으는 network 내부의 normalized slot. 물리 vertex나 외부 anchor graph가 아니다.\\
Response package & Learned area와 homogeneous rule로 파생한 mass measure, Global/Local motion field,
angular-frequency pole와 modal damping ratio, selector metadata를 포함한 setup output.\\
Dynamic state & Runtime generalized coordinate와 velocity. Wind load와 response package로 deterministic하게 갱신한다.\\
\bottomrule
\end{longtable}

\section{연구 질문, 가설과 기여 계층}

\subsection{연구 질문}

독립적으로 reconstruction한 static 3DGS asset \cite{kerbl2023gaussian}과 명시적인
material, metric scale 및 attachment만으로, target mesh나 target motion video 없이
바람에 대한 Global sway와 필요한 Local flutter response를 안정적으로 생성할 수 있는가?

핵심 가설은 세 부분이다.

\begin{enumerate}
  \item Mesh topology를 inference output으로 복원하지 않아도, training-only privileged relation과
  multi-resplat consistency가 patch-to-token encoder에 same-surface 구조를 내재화할 수 있다.
  \item Absolute next deformation보다 mass-normalized passive response package를 학습하면
  wind, material과 attachment 변화에 더 controllable하고 장기 rollout이 더 안정적이다.
  \item Strong fixed-rank Global package가 놓치는 spatially localized teacher response가 존재한다면,
  그 residual만 Local slot으로 학습하고 필요할 때만 실행하여 same-p95 quality--latency Pareto를 개선할 수 있다.
\end{enumerate}

\subsection{검증 대상 기여}

\textbf{검증 후 방어할 핵심 주장:} 서로 다르게 재구성된 static GS로부터 면적·질량과 바람의 일에
일관된 response package를 예측하고, 현재 바람에 필요한 국소 반응만 실행하여 지원 범위의 새로운
물체에 target mesh나 대상별 재학습 없이 적용한다. 아래 MC1--MC3는 이를 구성하는 검증 대상이며,
설계 또는 teacher 개발 진단만으로 기여가 입증되었다고 쓰지 않는다.

\begin{enumerate}
  \item \textbf{MC1 --- GS-measure-consistent latent response representation.}
  같은 물체의 independent GS density/seed 변화에도 표면 measure, total wind work와 response가 일관되도록
  patch 기반 latent token과 per-Gaussian area ownership 및 homogeneous mass quadrature를 response supervision과 함께 학습한다.
  여러 밀도의 GS를 만드는 것은 학습 수단이며, 기여는 표현 차이로 인한 가짜 질량·외력·응답 변화를
  억제하는 학습과 구성 방식이다. 독립 재구성 검증과 별도로 split--merge stress test를 둔다.
  \item \textbf{MC2 --- topology-privileged, target-mesh-free response distillation.}
  Training mesh의 vertex, relation과 simulation trajectory를 privileged teacher로 쓰되,
  inference에서는 새로운 static GS와 명시적인 material/metric/attachment로부터 passive response package를
  한 번 예측하여 cache한다. Target mesh, target motion video 및 per-target physical fitting이나 재학습은 요구하지 않는다.
  Mode column 자체나 다음 위치만을 학습하는 대신 force/impulse/frequency probe의 physical response를 맞춘다.
  이후 바람 입력에 대한 운동은 이 package를 시간 전진시켜 생성한다.
  \item \textbf{MC3 --- wind-residual-driven selective Global--Local response.}
  항상 활성인 Global package와 Global이 설명하지 못한 localized Local residual dictionary를 분리하고,
  fixed Local budget에서 현재 wind가 실제로 excite하는 slot만 실행한다.
  분해 자체보다, 더 큰 Global 및 always-on Local과 비교한 정확도--실행 시간 이득이 기여의 근거다.
  \item \textbf{SYS --- renderer-ready Gaussian response execution.}
  Area-weighted wind coupling, exact attachment, deterministic stable update와 Gaussian mean/covariance transport를
  하나의 측정 가능한 setup/runtime system으로 닫는다. 이는 MC1--MC3의 실행과 검증을 완성하는
  시스템 항목이며, 각 구성 요소의 최초성을 별도 주요 기여로 세지 않는다.
\end{enumerate}

\paragraph{Teacher의 역할과 기여의 경계.}
해상도에 따른 수치적 차이가 학습 신호를 지배하지 않도록 teacher의 공간·시간 수렴을 확인하는 것은
데이터 품질의 전제다. 기존 thin-shell 방법의 구현과 여러 GS에 대한 정답 mapping만으로 새로운 AI
방법론을 주장하지 않는다. 현행 P3 teacher는 Noels (2009)의 DG shell 요소 연결 원리를 참고한
Koiter/StVK 구현이며, 원 논문 전체의 재현이나 개선 및 그 안정성 보장의 승계를 주장하지 않는다
(상세 관계와 검증 범위는 \href{development/r1_teacher_probe_oracle.pdf}{R1의 P3 유한 회전 shell 절}).
Student의 기여는 특정 자체 솔버를 사용하는 데 의존하지 않는다. 동일한 teacher 계약과 정확도를
충족하는 기존 솔버도 원칙적으로 사용할 수 있으며, 실제 교체는 별도 검증과 채택을 거친다.

\paragraph{기여를 뒷받침하는 증거.}
MC1은 단순한 밀도 증강보다 measure/response consistency의 효과가 있어야 하고,
MC2는 같은 물체의 재표현을 넘어 object-disjoint 대상에서 검증한다.
MC3는 Local residual이 teacher/mapping 오차보다 크다는 전제와 selector 비용을 포함한
Global/always-on Local 대조를 요구한다. 해당 기준과 실패 시 claim 축소는 기존 Gate A--D를 따른다.
Teacher 수렴 통과만으로 학습 아티팩트 감소를 주장하지 않으며, 실제 학습 결과로 확인한다.

Latent token, attention, damped oscillator, virtual-work projection, covariance rotation 각각의 최초성을 주장하지 않는다.
Simplicits와 FreeForm은 mesh-free reduced simulation과 Gaussian splat 적용을 이미 다루며
\cite{modi2024simplicits,xiang2026freeform}, NeuROK과 LaGSplat은 learned latent kinematics/dynamics와
force pullback을 다룬다 \cite{geng2026neurok,pottier2026lagsplat}.
방어할 교차점은 \textbf{static independent GS에서 amortized wind-response prediction},
\textbf{GS measure/density consistency}, 그리고 \textbf{selective Local residual execution}이다.

\subsection{명시적 비주장}

\begin{itemize}
  \item Fully meshless training, relation-free network 또는 arbitrary topology recovery
  \item Unknown material, thickness, metric scale와 attachment를 appearance만으로 식별
  \item Engineering-grade stress, support reaction, exact shell constitutive behavior
  \item Wake, vortex shedding, aerodynamic history와 two-way fluid--structure interaction
  \item Self-contact, severe folding, tearing, plastic crease, free flight와 topology change
  \item 모든 물체 family에 대한 universal zero-shot dynamics
  \item Free-response passivity만으로 arbitrary time/state-dependent wind에서도 무조건 bounded하다는 보장
  \item Profiler와 matched-p95 evidence 전의 real-time 또는 speedup claim
\end{itemize}

정확한 표현은 \textbf{target-mesh-free inference without a persistent physical adjacency graph}다.
Encoder 내부 patch, attention affinity와 transient neighborhood까지 없다는 뜻은 아니다.

\section{범위와 I/O contract}

\subsection{Setup 및 frame 입력}

\begin{equation}
  \mathcal I_{\mathrm{setup}}
  =\left(\mathcal G^0,\ \vartheta_{\mathrm{metric}},\ \vartheta_{\mathrm{mat}},\
  \mathcal C_{\mathrm{attach}},\ \vartheta_{\mathrm{aero}},\ \vartheta_{\mathrm{domain}}\right),
  \label{eq:rd-setup-input}
\end{equation}
여기서 \(\mathcal G^0=\{\mu_i^0,\Sigma_i^0,\alpha_i,c_i\}_{i=1}^{N}\)는 rest static GS다.
\(\vartheta_{\mathrm{metric}}=(L_0,A_{\mathrm{ref}},M_{\mathrm{ref}})\)는 target에서 제공하는
positive world-length/area/mass reference이며, \(\vartheta_{\mathrm{mat}}\)는 stiffness/modal-damping preset,
\(\mathcal C_{\mathrm{attach}}\)는 hard attachment mask 또는 geometric region이다.
\(\vartheta_{\mathrm{aero}}\)는 domain 전체에 고정하는 effective traction scale
\(\kappa_{\mathrm{aero}}\)와 \code{traction_law_id}를 소유하고,
\(\vartheta_{\mathrm{domain}}\)은 허용 geometry, wind, timestep/frequency, numerical traction guard 및
OOD envelope를 소유한다. 이 typed physical metadata를 모르는 target은 core input 밖이며 appearance에서 역추론하지 않는다.
% [논문 작성 참고 -- 수정 이유]
% 현재 one-air-condition 범위에서는 공기밀도와 법선 저항계수가 곱으로만 힘에 나타나므로 각각 식별할 수 없다.
% 따라서 rho_air와 C_n을 별도 사용자/학습 파라미터로 두지 않고 kappa_aero=0.5*rho_air*C_n으로 합쳤다.
% Wind Projection Basis도 rho와 C_D를 하나의 scalar factor로 합치는 선례를 제공한다.
% 논문 citation 후보: \cite{diener2009windbasis}.
Appearance는 renderer가 소유하며 physics encoder는 opacity, covariance와 low-order appearance cue 중
development에서 허용한 subset만 사용한다. Texture memorization을 막기 위해 appearance randomization ablation을 둔다.

V0의 외부 physical contract는 nondimensional unit mode를 허용하지 않는 \textbf{SI-only}다.
즉 \(L_0\,[\mathrm m]\), \(A_{\mathrm{ref}}\,[\mathrm m^2]\), \(M_{\mathrm{ref}}\,[\mathrm{kg}]\),
\(\Delta t\,[\mathrm s]\), force \([\mathrm N]\)와 traction \([\mathrm{N/m^2}]\)를 사용한다.
\(\Omega_S\)는 angular-frequency convention의 diagonal matrix \([\mathrm{rad/s}]\)이고
\(Z_S=\operatorname{diag}(\zeta_{S,j})\)는 dimensionless modal damping-ratio matrix다.
\(M_{\mathrm{ref}}\)만 total mass를 소유하며 material preset이나 network output이 두 번째 total-mass
scale을 만들 수 없다. Unit mode, angular-frequency convention과 typed scalar의 단위는 package/run hash에 기록한다.
Encoder의 \(\widehat\mu,\widehat\Sigma\)는 numerical feature normalization일 뿐 별도 physical unit mode가 아니며,
response loss, package validation과 evaluator에 들어가기 전 모든 physical payload를 위 SI convention으로 복원한다.

Frame 입력은
\begin{equation}
  \begin{aligned}
  \mathcal I_t&=\left(W_t(\cdot),\ \vartheta_{\mathrm{ctrl},t},\ \Delta t,\ s_{t-1},\ \mathcal P\right),\\
  \vartheta_{\mathrm{ctrl},t}&=\left(g_W(t),g_K,g_D\right),
  \end{aligned}
  \label{eq:rd-frame-input}
\end{equation}
이다. \(W_t\)는 prescribed wind field, \(s_{t-1}\)는 cached generalized state,
\(\mathcal P\)는 setup response package다. Package는 rest GS와 모든 target-allowed physical metadata의
identity/hash 및 runtime payload를 소유하므로 frame API가 teacher mesh나 별도 hidden state를 조회하지 않는다.
Target dynamic RGB/video나 target mesh state를 재입력하지 않는다.

\subsection{Typed artist control의 물리적 의미}

Artist control은 자유로운 latent vector \(\mathbf a\)로 network에 추가하지 않고,
이미 선언한 typed input과 cached response의 유효 범위 안에서만 정의한다.
Development에서 동결한 \(g_W(t)\in[0,g_W^+]\)는 바람의 크기만
\(W_t^{\mathrm{eff}}=g_W(t)W_t\)로 조절하고 방향은 \(W_t\)가 소유한다. Package의
\(\Omega_S^{\mathrm{base}},Z_S^{\mathrm{base}}\)는 immutable이고,
\(S\in\{G\}\cup\{(L,k)\}_{k=1}^{K_L}\)에 대한 evaluator가
\begin{equation}
  \Omega_S^{\mathrm{eff}}=\sqrt{g_K}\,\Omega_S^{\mathrm{base}},
  \qquad
  Z_S^{\mathrm{eff}}=g_DZ_S^{\mathrm{base}},
  \qquad g_K>0,\ g_D>0
  \label{eq:rd-typed-control-map}
\end{equation}
를 파생한다. \(g_K,g_D\)는 package에 저장된 허용 구간 안에서 rollout 전체에 고정하고,
cached base field를 in-place로 수정하지 않는다. 둘 중 하나를 바꾸면 derived evaluator identity/hash를
갱신하고 rest state에서 새 rollout을 시작한다. 이 조절은 arbitrary material을 추론하는 기능이
아니며, material preset 자체나 attachment를 바꾸면 setup network를 다시 실행해 새 package를
만들고 rest state로 reset한다. Run manifest는 rollout의 actual \(g_K,g_D\)와 \(g_W(t)\)
program/hash를, frame trace는 각 frame의 actual \(g_W(t),g_K,g_D\)와 derived evaluator identity를 소유한다.

각 control은 동일한 rest state와 wind program에서 amplitude, dominant frequency, phase와
free-decay envelope의 반응을 sweep한다. Wind gain 대비 load/response direction, stiffness gain 대비
frequency ordering, damping gain 대비 decay ordering과 in-envelope 유한성을 보고, expected monotonicity나
phase continuity가 깨지면 해당 control 범위를 축소하거나 비주장한다. Optional Local/flutter style
gain은 Gate C 통과 후에만 Global을 건드리지 않는 정성 demo control로 분리하며,
Gate B/C 정량 evidence에 사용하지 않는다.

\subsection{정량 core와 champion domain}

\begin{longtable}{L{0.21\linewidth}L{0.34\linewidth}L{0.35\linewidth}}
\toprule
축 & 정량 core & Gate 후 확장\\
\midrule
\endfirsthead
\toprule
축 & 정량 core & Gate 후 확장\\
\midrule
\endhead
Geometry & flat-rest single-layer strip, rectangular flag & triangular pennant, simple curtain panel, mild curved silhouette\\
Attachment & 한 edge 또는 작은 declared patch의 hard clamp & dev에 포함한 제한적 attachment family\\
Material & 2--3 known isotropic presets & preset interpolation; arbitrary inference 금지\\
Wind & steady, start/stop, sinusoidal, spatial/traveling gust & unseen composition과 제한적 field perturbation\\
Motion & small-to-moderate deformation, one-way wind & 더 큰 rotation은 confidence/OOD와 함께 qualitative\\
Rendering & mean + covariance rotation, degree-0 physics evaluation & higher SH/frame transport는 paper-stage\\
\bottomrule
\end{longtable}

Champion scene은 사후에 성공 사례만 고르는 뜻이 아니다. 위 domain과 failure criteria를 development에서 동결하고,
그 안에서 Global sway, free-edge flutter와 traveling gust가 잘 드러나는 3--5개 hero scene을 선택한다.
같은 범위의 object-disjoint 작은 quantitative suite와 실패 장면도 함께 공개한다.

\section{Training-only teacher와 paired 3DGS data}

\subsection{Teacher가 소유하는 정보}

Teacher mesh simulator는 vertex/element connectivity, surface quadrature, material, attachment와
prescribed wind를 소유한다. Reference trajectory에는 canonical surface sample에서의
position, displacement, velocity, optional normal, applied traction와 external work를 저장한다.
Teacher refinement와 timestep convergence가 확인되기 전에는 student training을 시작하지 않는다.
Core teacher와 target GS evaluator는 각자의 surface quadrature 위에서 같은
\code{traction_law_id}, effective scale, guard-before-area 순서와 frame-start sampling contract를 사용한다.
다른 teacher forcing law는 V0 label로 섞지 않고 별도 dataset/method version으로 분리한다.
% [논문 작성 참고 -- 수정 이유]
% tau_guard는 재료나 공력의 물리계수가 아니라 수치 안전장치다. 동일 teacher/student identity에는 포함하되
% coefficient tuple로 해석하지 않고 정상 quantitative envelope에서 activation zero를 요구한다.

Mesh vertex를 모두 supervision에 사용하는 것은 유용하지만, \textbf{vertex 하나를 latent anchor 하나와
일대일 대응시키지 않는다}. Vertex 수와 tessellation은 simulator discretization 선택이기 때문이다.
선형 nodal teacher에서는 vertex와 area-weighted sample을 연결할 수 있지만, 일반적인 고차 teacher의
자유도는 양의 면적을 가진 물리 sample과 같지 않다. Teacher field를 양의 quadrature와 canonical
surface probe에서 평가한 뒤 fixed/capped latent token의 response supervision에 사용한다.
공통 probe로의 remap은 서로 다른 teacher mesh 및 GS representation의 정량 비교에서 필수다.

Teacher의 shape-evaluation matrix를 \(\Psi\), 양의 quadrature를
\(W_{A,T}=\operatorname{diag}(a_q)\), \(\sum_q a_q=A_{\mathrm{ref}}\)라 하면,
homogeneous consistent mass를 사용하는 후보는
\begin{equation}
  M_T^{\mathrm{scalar}}=\frac{M_{\mathrm{ref}}}{A_{\mathrm{ref}}}
  \Psi^\top W_{A,T}\Psi,\qquad
  \mathbf 1^\top M_T^{\mathrm{scalar}}\mathbf 1=M_{\mathrm{ref}},\qquad
  M_T=M_T^{\mathrm{scalar}}\otimes I_3
  \label{eq:rd-teacher-consistent-mass}
\end{equation}
를 따른다. 총질량 합 식은 고정부 DOF를 제거하기 전의 전체 DOF 질량행렬에 대한 조건이다.
고정부를 제거한 행렬에 같은 총합을 요구하지 않는다. 질량행렬의 양정성은 충분한 quadrature rank와
허용 자유도에서 검증한다.
Teacher element/quadrature/mass law를 registry에 명시하며, 고차 자유도에 임의의 양의 point mass를
붙이지 않는다. 이 확장은 student GS의 기존 positive area와 diagonal mass 계약을 바꾸지 않는다.

\subsection{Mandatory common-probe transfer contract}

Training loss와 모든 paired quantitative evaluation은 source object마다 test open 전에 고정한
canonical probe set
\[
  \mathcal X_P^0=\{(p_j^0,w_{A,j}^P,w_{M,j}^P)\}_{j=1}^{N_P}
\]
에서 수행한다. Teacher discretization을 \(X=T\), independent GS variant를 \(X=r\)라 할 때,
ordered interpolation support와 representation별 map law를 갖는
\begin{equation}
  T_X\in\mathbb R^{N_P\times N_X},\qquad
  \sum_i(T_X)_{ji}=1,\qquad
  S_X=T_X\otimes I_3
  \label{eq:rd-common-probe-map}
\end{equation}
를 만든다. GS의 convex support map과 선형 barycentric teacher map에는
\((T_X)_{ji}\ge0\)를 유지한다. 고차 teacher shape function에만 부호가 있는 weight를 허용하며,
별도 map-law/version으로 polynomial reproduction, conditioning/noise amplification과 coverage를
검증한다. Signed shape weight와 양의 surface quadrature는 서로 다른 물리량이다.
기존 nonnegative map schema를 암묵적으로 재해석하지 않는다. Reference와 GS response는
\[
  u_P^*=S_Tu_T^*,\quad v_P^*=S_Tv_T^*,\qquad
  u_P^{(r)}=S_ru_{\mathrm{GS}}^{(r)},\quad
  v_P^{(r)}=S_rv_{\mathrm{GS}}^{(r)}
\]
로 같은 공간에 올린다. 모든 accepted \(S_X\)는 row-partition tolerance와 임의의
\(d\in\mathbb R^3\)에 대한 constant-field reproduction
\[
  S_X(\mathbf 1_{N_X}\otimes d)=\mathbf 1_{N_P}\otimes d
\]
를 통과해야 한다. 추가 mapping-noise 기준은 source rest position에서 평가한 predeclared
constant/coordinate-linear analytic vector-field family를 probe로 보낸 normalized reproduction error로
측정하고, 모든 \(X\)가 dev-frozen tolerance를 통과해야 한다. Manifest는 probe ID/rest position/area--mass weight, 각 row의 ordered
support와 weight, coverage distance, unsupported-probe rate, kernel/version, \(S_T,S_r\) hash,
fixture tolerance와 측정된 mapping-noise floor를 저장한다. Coverage가 dev-frozen threshold보다
낮은 variant는 결과를 본 뒤 probe를 골라내지 않고 predeclared failure/reject 규칙을 따른다.
정량 denominator는 teacher와 source-object group의 모든 declared GS variant가 함께 support하는
probe intersection으로 만든 하나의 pre-test frozen common-valid mask를 사용한다. Variant마다 유리한
probe subset을 따로 선택하지 않으며 mask, excluded reason과 hash를 모든 paired run이 공유한다.

Adjoint는 payload의 물리량을 구분한다. \(f_P\)가 probe별 \emph{total force}이면
\(f_X=S_X^\top f_P\)가 virtual work를 보존한다. 반면 \(\tau_P\)가 probe
\emph{traction}이면 \(W_{A,P}=\operatorname{diag}(w_{A,j}^PI_3)\)를 먼저 적용해
\begin{equation}
  f_X=S_X^\top W_{A,P}\tau_P,\qquad
  \tau_P^\top W_{A,P}S_Xu_X=f_X^\top u_X
  \label{eq:rd-common-probe-adjoint}
\end{equation}
인 source total force를 만든다. Traction에 \(S_X^\top\)만 직접 적용하거나 area를 두 번
곱하는 구현은 금지한다. 이 transfer는 training/evaluation fixture 전용이다. Target runtime의
GS traction과 force는 common-probe map 없이 각 Gaussian에서 직접 계산한다.

Teacher convergence도 두 축을 섞지 않는다. Accepted fine timestep을 고정한 spatial mesh-refinement
error와 accepted fine mesh를 고정한 temporal timestep-refinement error를 동일 probe와 physical-time
grid에서 별도로 측정하고, 각각 dev-frozen tolerance와 hash를 통과해야 한다. 두 refinement를 동시에
바꾼 차이만으로 cancellation을 허용하지 않으며, 이 mapping/convergence fixture가 끝나기 전에는
student response fitting을 시작하지 않는다.
두 축의 acceptance는 같은 structure/element/mass/damping/attachment/load identity를 공유해야 한다.
서로 다른 teacher 후보의 시간 통과와 공간 통과를 합쳐 하나의 reference를 승인하지 않는다.
정적 polynomial energy 일치만으로 동적 수렴을 대신하지 않으며, interior weak-force consistency,
mesh 방향, 초기 상태의 고주파 성분과 velocity 오차도 진단한다.

\subsection{Teacher 개발 근거와 현재 채택 경계 (2026-09-09)}
\label{sec:rd-teacher-development-evidence}

현재 구현에는 Newton native trajectory/registry/probe 경로, CPU nonlinear shell 진단,
P2/P3 \(C^0\) interior-penalty 선형 판 및 독립 \(C^2\) cubic B-spline 기준이 있다.
고정 native hinge 계수의 mesh 의존성, 면적 보정 hinge의 방향 편향, quadratic patch의 내부 힘
결함을 확인했다. 따라서 이 후보들의 정적 에너지 검사를 곧바로 accepted teacher 근거로 쓰지 않는다.
P3 후보는 고정된 처방 압력 fixture에서 공간/방향/독립 기준의 1\% 개발 기준을 통과했지만,
원래 \(u_z^0=10^{-3}x^2\) 초기 변위의 속도 수렴은 실패 상태다.
여기서 \(x=0\) 경계는 위치만 고정하고 회전 기울기는 자유다. 이를 clamped 판으로 소개하지 않는다.

앞선 P3 pressure 결과는 선형 normal plate의 해당 하중에 한정된다. 아래의 작은 굽힘 wind/reset은
새 경계·상대풍·frame-start hold를 별도로 검증했다. Full nonlinear membrane+bending,
새 고차 public probe/registry, 해당 nonlinear backend의 시간·공간 검증 및 independent-GS oracle는
아직 연결되지 않았다. 기존 CPU shell의 시간 수렴 통과를 P3에
승계하지 않는다. Development loader로 검증한 15개 Newton sample window는
\code{training_eligible=false}를 유지하며 R2 학습 label이 아니다.
1\%는 해당 개발 fixture의 사전 고정 기준이지 모든 R1 metric에 새로 적용한 보편 threshold가 아니다.

전체 설계·반례·수치·재현 경로는 \code{development/r1_teacher_probe_oracle.tex}의 관련 절에 통합한다.
해당 문서의 본문 수식을 직접 갱신하고 변경 이유와 적용 범위를 인접 설명·TeX 주석으로 보존한다.
실험용 처방 압력과 rest-start 입력은 원래 displaced free-decay 요구를 삭제하지 않는다.
첫 accepted 데이터의 입력 범위는 아직 미동결이다. Rest에서 시간 변화하는 랜덤 바람으로
연속 궤적을 만들고, 중간 형상에서 속도만 0으로 만든 뒤 계속하는 개발 비교를 실행했다.
구체적인 절차와 판정은 R1의 초기 상태 탐색 절이 소유한다.

\subsection{초기 상태 분포와 속도 초기화 탐색}
\label{sec:rd-initial-state-exploration}

최종 supported envelope에서는 이미 변형되고 움직이는 상태에서도 새 바람에 반응해야 한다.
Rest-start 궤적도 중간에는 이러한 상태를 포함한다. 변형 상태에서 시작하는 데이터의 추가 가치는
rest-start 바람으로 드물게 방문하는 형상·속도 조합을 보충하는 데 있으며, 학습 후 개선은 미검증이다.
동일한 rest/material/attachment와 완전한 물리 상태 및 이후 바람이 같으면, 시작 경로만으로 다른
정답 운동을 정의하지 않는다. 형상만 같고 속도가 다르면 서로 다른 초기 상태다.

이 비교에서는 전체 물체의 rest 형상을 유지한다. 선택한 시점의 원본 궤적에서
자연스럽게 계속하는 분기와, 위치는 보존하고 전체 자유 DOF의 속도를 0으로 만드는 분기를 만든다.
두 분기는 같은 절대 시각의 이후 wind field를 사용한다. Relative-wind traction은 각 분기의
현재 속도에서 다시 평가하므로 applied force까지 같다고 가정하지 않는다.
속도 초기화는 외부 개입에 의한 운동에너지 제거이며 자연 감쇠나 공력 off가 아니다.
2026-09-09 native Newton 개발 경로에서 5개 조건의 원본/독립 replay/reset 25개 trace를 생성·검산했다.
재생의 위치/속도/힘/work 차이는 0이고 개입 시 위치 보존과 운동에너지 제거를 확인했다.
하지만 자연 연속의 finest 공간/시간 속도 차이는 각각 12.0401\%/2.1767\%로 개발 기준 1\%를 넘었다.
1.1 s reset의 값은 81.1482\%/17.7560\%이며 작은 fine 속도 분모와 절대 차이를 함께 해석해야 한다.
현재 결과는 frame 표본의 진단이며 accepted Teacher/data scope 또는 학습 이득이 아니다.

또한 모서리의 위치만 고정한 조건은 힌지처럼 강체 회전할 수 있다.
0.7 s 형상은 고정 모서리 주위 회전을 제거하면 probe RMS 잔차가 0.002380 mm로,
전체 위치 변화 RMS 12.700918 mm에 비해 작았다. 큰 위치 변화가 큰 내부 굽힘 상태를 뜻하지 않는다.
후속 비교에서 사용자가 선택한 왼쪽0.25 m 영역 고정으로 비평면 굽힘을 확보했다.
그러나 자연 연속의 공간8$\to$16 속도 차이는140.2337\%였고, 시간 보완64$\to$128 후에도
네 분기 모두 변위/속도1\%의 동시 기준을 통과하지 못했다.
같은 고정 조건의 해석적 굽힘 형상에서 native 에너지는 mesh4$\to$32에8.888857$\to$1.435179 $\mu$J로 감소했다.
이는 시간 적분 없이 재현되는 물성의 격자 의존성이다. 사용자는 \textbf{기존 P3의 작은 굽힘 범위 활용}을 선택했다.
Native 모델을 고쳤다고 주장하지 않고, 별도 P3 자유 판과 고정 경계·consistent mass에서 다시 검증했다.
정량 근거와 시간/반복 분리, 원본 및 재현 경로는 R1에 통합했다.
% [논문 작성 참고 -- 결정 이유]
% 2026-09-09 사용자 결정: rest에서 바람으로 도달한 변형을 사용해 속도만 0으로 만든 후 반응을 비교한다.
% 임의 초기 변형의 무제한 허용, 기존 x^2 실패 삭제 또는 rest 형상 재정의로 해석하지 않는다.

후속 P3 조건은 \(E=10^6\) Pa, \(\nu=0.3\), \(h=0.01\) m와 면밀도0.1 kg/m$^2$다.
기존 seed/60 Hz/1.5 s/reset 시점은 유지하고 target 풍속 상한을0.05 m/s로 낮췄다.
자유 길이0.75 m의 변위·기울기 고정, current relative-wind traction과 전체 모드의 exact held-force 전진을 사용한다.
P3 공간8$\to$16의 네 분기 중 최악 속도 오차 상한은0.882329\%, 독립 spline 비교는0.386575\%로
고정1\% 기준을 통과했다. 요소 내부와 모든 시간의 변위 상한0.421778 mm 미만으로 \(0.1h=1\) mm를 지켰다.
독립 재실행의 전체 배열과 에너지/work/reset 검산 뒤 \textbf{19 window/76 patch sample을 생성·검증}했다.
이는 검증된 작은 굽힘 fixture의 적격성이다. 큰 변형·새 object/seed·장기 안정성·학습 이득이나
R1 전체 완료를 뜻하지 않는다. 기존 P3 pressure, native 실패와 이번 새 경계/공력의 증거를 구분한다.

\textbf{후속 진행 원칙 (2026-09-09):} 사용자는 물리 수식을 제대로 구현·검증한 다음에
학습데이터를 생성하도록 지시했다. 기존 76개 sample은 역사적 개발 근거로 보존하고 추가 발행은
보류한다. P3의 3D 요소 내부 기하, 질량, 상대풍 및 응력 가상 일의 기하 미분을 보완했다.
후속 사용자 결정은 \textbf{1번 유한 회전 shell}이다. 기존 P3에 현재 metric의 Green strain,
현재 법선 곡률과 Koiter/StVK 막·굽힘 저장 에너지를 연결했다. 이는 작은 재료 strain을 전제하며
큰 회전을 허용하는 고정 두께 thin-shell 범위다. 요소 연결은 normal jump와 current moment flux의
보존적 edge 에너지, 고정 경계는 full flux와 fixed normal을 사용한다.
전체 에너지의 힘/HVP, consistent mass와 비선형 Newmark를 함께 구현했다.
정적 원통36개 조합의 finest 에너지 오차는0.007472\% 미만이고,
rest/굽은 초기 상태의 짧은 시간 응답은 독립 DOP853 기준과 일치했다.
정적 수식과 짧은 시간 대조를 큰 변형 Teacher의 동적 공간 수렴·장기 안정성으로 승계하지 않는다.
원통을 지지 하중 없이 놓는 진단에서는 자유단 moment 조건의 불일치와 강한 고주파 응답을 확인했다.
시간 세분화로 n4의 속도 차이를0.7622\%로 줄였지만 거친 메시의 공간 차이는 크게 남았다.
이 진단을 바람으로 도달한 상태와 구분하고, 사용자가 정한 rest-start/velocity-reset 절차를 유지한다.
별도5 m/s peak의 짧은 수식 진단에서 약7.3 cm 변형과0.254 rad 회전에 도달한 뒤 reset/계속 전진을 검산했다.
이는 데이터 풍속 범위의 동결이나 큰 변형의 공간 수렴 완료를 뜻하지 않는다.
약한 바람3 frame의 공간 n16$\to$32 속도 차이는0.385696\%, n16의 대각선 방향 차이는0.153774\%다.
5 m/s reset의 초기 n8 시간 비교는0.304620\%, 공간 n8$\to$16은1.849529\%였다.
공간 기준 실패를 보존하고 아래 GPU 세분 비교에서 n16$\to$32의 통과를 확인했다.
P3 공간과 Newmark의2차 시간 보간에 대한 Bernstein 계수로 요소 내부·시간 간격 사이의
기하·strain 상한을 추가했다. 이 수치 보간의 단사성 충분조건과 정확한 연속 운동의 보증은 구분한다.
상세 수식, 고정 법선 반력 보완, 실패·수렴 수치와 원본 경로는 R1의 P3 유한 회전 shell 절에서 관리한다.
후속 사용자 선택에 따라 동일 에너지·힘·정확한 방향 HVP와 현재 면적 공력을 Warp float64로 이식했다.
실제 GTX1080Ti의12개 상태 연산자 검사와 n16의3 frame/reset CPU 대조를 통과했다.
Newton 제어와 희소 선형 풀이는 CPU에 유지한다. GPU 구현 일치는 물리 수렴의 대체 조건이 아니다.
0.05 s의5 m/s peak/reset 진단에서는 n16$\to$32/sub256 공간 속도 상한0.400065\%,
n32/sub128$\to$256 시간 상한0.243747\%, n32/sub256 방향 상한0.040619\%를 얻었다.
모두 수치 보간 전체 시간의1\% 기준을 통과했고 원식은 CPU로 별도 검산했다.
이후1.5 s 검증은 사용자가 선택한 기존0.25--0.5 m/s knot로 수행했다.
Rest natural과0.3/0.7/1.1 s의 변형 상태에서 velocity를0으로 둔 독립 분기를 비교한다.
동일 전역 시각의 바람을 유지하며, 완료된 parent checkpoint를 참조해 prefix 반복 계산을 줄인다.
세부 저장·검산·수렴 결과는 R1에 누적하고 추가 학습데이터 발행은 보류한다.
다섯 공간·시간·방향 조건의 자연 응답과 세 reset, 총20개 원본/239,616 interval을 검산했다.
다섯 비교의 네 분기20개 모두 전체 수치 보간의 변위·증분·속도1\%를 통과했다.
공간 n8$\to$16의 최대 속도 상한은 sub128의0.910953\%에서 sub256의0.904299\%로 유지됐고,
시간 sub128$\to$256의 최대 상한은 n8의0.169299\%, n16의0.148473\%였다.
N16/sub256의 방향 최대 상한은0.063805\%다.
CUDA 전 interval 검산과 frame별 CPU3상태/공력 대조, checkpoint 재시작 및 선택5상태의 구적 대조를
구분해 R1에 보존했다. 약9 mm 변위에서 velocity-reset/동일 미래 바람 전진의 수치적 성립을 확인했지만,
큰 변형의 장시간 coverage와 독립 비선형 공간 기준·R1 전체 채택은 미완료다.
후속 사용자 승인 범위는 동일 파형4배(knot1--2 m/s)의 자연 응답·세 reset과 원식·공간·시간·방향
검증까지다. 물리식과 허용오차를 유지하며, 풍속의 추가 확대·학습 적격성 채택·데이터 발행은 제외한다.
4배 입력의 n16/sub256 natural은 원식·기하 검산을 통과했고 최대 nodal 변위는109.084864 mm였다.
그러나 n8$\to$16/sub128 공간 속도 상한은 natural/0.3 s/0.7 s/1.1 s reset에서 각각
1.087962/1.211015/1.193883/3.495261\%로 기준1\%를 넘었다. 이 실패를 보존하고 n16$\to$32의
공간 비교 및 n32의 시간·방향 비교를 추가했다. N16 natural의 시간·방향 속도 상한은
각각0.144732/0.059589\%로 통과했다. 이어서 세 reset까지 모두 확인한 n16 시간·방향 비교의
네 분기 최댓값도 각각0.731830/0.274527\%로 통과했다. N32/sub128 natural의 원식 검산과
정확한 재시작도 통과했다. N32/sub256 natural의23,040 interval 원식·기하 검산과
n16$\to$32/sub256 공간 속도0.444132\%, n32 시간 속도0.140173\% 비교도 통과했다.
최대 nodal 변위는109.091615 mm다. 후속 저장 요약에서는 n32 시간·방향의 네 분기는 모두 통과하나
n16$\to$32/sub256의1.1 s reset 공간 속도 상한1.855258\%가1\% 기준을 넘는다.
필수20개 비교 중19개 통과이며,2026-09-11 report/summary 대조와 전체 원본/물리 재검산의 범위는 구분한다.
Coarse 공간 비교의 sub128과 보완 비교의 sub256을 구분하고, 두 값만으로 공간 수렴 차수를 추정하지 않는다.
Wind 배율·parent·검산 identity 및 기존/개선 producer 구간의 보존 계약은 R1에 기록한다.
4배 계산은 종료했지만 최종 수렴 기준 미달이며 R1 채택과 추가 학습데이터 발행은 보류다.
후속 Teacher 실행 경로는 CPU 반복 풀이+GPU HVP/CUDA graph를 채택했고, 전체1.5 s의 잠정 성능
비교와 CuPy 희소 풀이 미채택 이유를 R1에 연결했다. 물리 수식·정밀도·허용오차는 바꾸지 않았다.
장면별10--20 s를 향한 첫10 s \(\Delta t\) 탐색은60 Hz 공력 갱신을 유지하는 최대256배 범위다.
큰 간격의 풀이 실패와 기존 간격의 시간 한도로 추천을 확보하지 못했으며, 이 결과로 물리적 불안정이나
Student의30fps 가능 여부를 단정하지 않는다. 정확한 조건과 미검증 범위는 R1의10초 탐색 절이 소유한다.
수치 응답 비교에는 reset 좌극한과 fine 시각을 포함하며,2차 Newmark 보간의 전체 시간 오차 상한을
추가했다. 이는 저장 수치 궤적의 보간에 대한 결과이고 정확한 연속 ODE 해의 인증은 아니다.
추가 학습데이터 발행은 계속 보류한다.

시간 window와 patch supervision은 전체 물체의 coupled trajectory에서 추출한다.
Reset 불연속을 가로지르는 구간을 외부 개입 정보가 없는 일반 wind-response label로 섞지 않으며,
같은 원본의 자연 연속/재시작 분기는 같은 source group으로 분할한다.
현재 setup-once student의 상태 초기화는 별도 계약이다. Teacher의 변형·속도 데이터를 추가한 사실만으로
임의 변형 GS에서의 runtime 시작 기능을 확보했다고 주장하지 않는다.
Teacher 수렴, 현재 response package의 oracle 재현 능력과 unseen-wind 장기 실행을 확인한 뒤에만
데이터 확대와 초기 상태 범위 채택을 판단한다.

\subsection{Independent multi-resplat pairing}

\paragraph{고정하는 물리 정답.}
물체별로 선택한 하나의 teacher 해상도에서 동일 초기 상태·물성·고정 조건·바람·시간 간격의 응답을
생성하고, 이를 같은 rest surface의 여러 GS variant에 mapping한다. 여기서 단일 해상도는 모든
물체에 같은 vertex 수를 강제한다는 뜻이 아니다. 생산 해상도의 타당성을 확인하는 별도 mesh/time
refinement는 유지한다. GS 밀도마다 teacher를 다시 풀어 서로 다른 정답을 만들지 않는다.
GS마다 점 수와 위치가 다르므로 같은 길이의 label 배열을 복제하지 않고, 동일 teacher 응답을
각 GS의 map으로 전달한 뒤 공통 surface probe와 area/mass metric에서 비교한다.

\paragraph{독립적인 GS 생성과 목적.}
같은 rest surface와 공통 이미지·카메라를 사용하고 GS 생성 시 Gaussian 예산 또는 densification 제한을
달리하여 각 GS를 독립적으로 최적화한다. 먼저 밀도 제한의 영향을 비교할 때는 initialization seed와
나머지 최적화 조건을 가능한 한 공통으로 두고, 이후 seed/camera subset 변화는 별도 축으로 검증한다.
Gaussian의 개수뿐 아니라 중심·크기·방향·opacity·겹침이 달라져도 같은 물리 응답을 예측하게 하는 것이 목적이다.
단순히 한 GS의 일부를 무작위 삭제하거나 split/duplicate한 것은 independent variant로 세지 않는다.
Random subsampling은 입력 누락에 대한 보조 증강 또는 비교군이며, 독립 재구성의 대체 근거가 아니다.
설정한 예산과 실제 Gaussian 수, 표면 복원 품질을 함께 기록한다. 지나치게 낮은 예산으로 형상 자체가
손실된 경우는 밀도 효과와 기하 정보 부족을 구분하며, 같은 밀도 일반화 판정으로 묶지 않는다.
같은 source object의 모든 mesh, GS variant, material과 wind sequence는 하나의 train/dev/test group에 둔다.
같은 물체의 다른 GS에 대한 일관성과 보지 못한 물체에 대한 일반화는 별도 결과로 보고한다.

권장 최소 단계는 다음과 같다.

\begin{enumerate}
  \item \textbf{K0 smoke:} strip 1개와 flag 1개, material 1개, independent GS 2개, impulse와 step wind.
  \item \textbf{K1 feasibility:} narrow family의 object-grouped train/dev/test, object당 2--3 GS density/seed variant,
  2--3 material preset과 네 probe family.
  \item \textbf{Paper evidence:} K1 Gate가 통과한 뒤에만 hero scene, real reconstruction과 broader silhouette를 추가한다.
\end{enumerate}

Core probe는 short impulse, step-on/off, log-chirp 또는 multi-sine, traveling Gaussian gust다.
Probe peak와 duration은 asset별로 조정하지 않고 development에서 domain-wide freeze한다.

\subsection{K0/R1 network-free oracle preflight}

Student network 전에 strip의 converged teacher와 independent GS 두 개로 common-probe oracle Global
fixture를 통과한다. 각 GS의 attached row를 exact zero로 둔 raw field를 teacher의 common-probe
impulse/FRF response에 fit하고, SI unit, force sign, \(S_T/S_r\), mass metric과 exact evaluator를
독립적으로 검사한다. 두 oracle package의 response를 다시 \(\mathcal X_P^0\)에 올린 paired error는
mapping과 evaluator contract의 preflight이지 MC1의 learned held-out evidence가 아니다.

Raw field \(B_c\)를 \(G=B_c^\top M_xB_c\), \(R=G^{-1/2}\),
\(B=B_cR\)로 mass-normalize하면 coordinate basis도 변한다. 따라서 teacher의 기존 diagonal
\(\Omega,Z\)를 그대로 붙이는 것은 금지한다. Oracle은 다음 둘 중 하나만 사용한다.
\begin{enumerate}
  \item normalized \(B\)를 고정한 뒤 common-probe time/FRF response로 positive
  \(\Omega,Z\)를 다시 fit한다.
  \item raw coordinate의 full reduced \(M_c,C_c,K_c\)를 보유할 때
  \(R^\top M_cR,R^\top C_cR,R^\top K_cR\)로 함께 변환한 뒤 일관되게
  re-diagonalize하여 field와 pole을 갱신한다.
\end{enumerate}
Arbitrary whitening과 unchanged diagonal pole의 조합, response fit 없이 mapped teacher mode만으로
oracle 성공을 선언하는 경로는 허용하지 않는다. Rank floor, attachment zero, mass orthogonality,
common-probe response와 free-decay exact-ZOH fixture가 모두 통과한 뒤에만 R2 network overfit으로 간다.

\subsection{Privilege boundary}

\begin{longtable}{L{0.29\linewidth}L{0.29\linewidth}L{0.31\linewidth}}
\toprule
정보 & Training & Target inference\\
\midrule
\endfirsthead
\toprule
정보 & Training & Target inference\\
\midrule
\endhead
Mesh vertex/connectivity & teacher simulation과 auxiliary label에 사용 & 금지\\
Mesh--GS correspondence & loss/evaluation remap에 사용 & 금지\\
Static GS attributes & student 입력 & 필수\\
Physical metadata & 명시 입력 & 명시 입력\\
Wind program & teacher force와 student rollout & runtime 입력\\
Dynamic target video & qualitative evaluation에만 선택적 & 방법 입력으로 금지\\
Persistent edge/triangle graph & teacher 내부에만 존재 & 생성·저장·소비 금지\\
\bottomrule
\end{longtable}

Forward API가 forbidden field를 받지 않는지, batch가 source-object split을 넘나들지 않는지와
checkpoint가 test normalization을 보지 않았는지를 manifest와 automated leakage test로 확인한다.

\section{Patch-to-token encoder와 내부 relation}

\subsection{Normalization과 Gaussian token}

Object center \(o\), positive length scale \(L_0\)와 attachment-defined canonical cue를 사용해
position/covariance를 normalize한다.
\begin{equation}
  \widehat\mu_i=\frac{\mu_i^0-o}{L_0},\qquad
  \widehat\Sigma_i=\frac{\Sigma_i^0}{L_0^2},\qquad
  g_i=E_{\mathrm{gs}}(\widehat\mu_i,\widehat\Sigma_i,\alpha_i,c_i,b_i,\vartheta_{\mathrm{mat}}).
  \label{eq:rd-gaussian-token}
\end{equation}
\(b_i\)는 attachment indicator 또는 distance cue다. Proper rigid transform과 scale 변화에 대한
canonicalization/equivariance fixture를 별도로 둔다. Severe symmetry에서 canonical frame이 모호하면
confidence reject 또는 sign-invariant feature를 사용한다.

\subsection{Patch 기반 latent anchor}

입력 GS를 normalized space의 overlapping local patch로 나누고, 각 patch evidence를 fixed/capped
\(K\)개의 latent query가 density-normalized cross-attention으로 읽는다 \cite{jaegle2021perceiverio}.
개념적으로
\begin{equation}
  h_k^{\mathrm{loc}}
  =\frac{\sum_i w_{ki}\,g_i}{\sum_i w_{ki}+\epsilon_w},\qquad
  w_{ki}\ge0,
  \label{eq:rd-latent-anchor}
\end{equation}
이며 실제 attention은 learnable key/query와 covariance/tangent cue를 함께 쓸 수 있다.
Patch는 계산 receptive field이지 물리 element가 아니다. \(h_k\)를 ``latent anchor''라고 부르더라도
그 좌표와 edge를 deployment package에 내보내지 않는다.

Patch-only로는 멀리 떨어진 attachment와 free edge의 관계를 놓칠 수 있으므로,
local aggregation 뒤 \(K\) token 간 global self-attention 또는 pooled object token을 사용한다.
즉 사용자가 원한 ``patch 기반이되 필요하면 전체 3DGS도 본다''를
\textbf{local evidence first, global context second}로 구현한다.

Token을 얻은 뒤에는 원래 Gaussian feature가 token과 object context를 다시 읽는 query decoder를 사용한다.
이 decoder는 input permutation에 equivariant하다.
\[
  y_i=D_{\mathrm{query}}\!\left(g_i,\{h_k\}_{k=1}^{K},h_{\mathrm{obj}}\right),
  \qquad i=1,\ldots,N .
\]
여기서 \(y_i\)는 Gaussian별 area, normal/validity와 Global/Local response-field row를 생성하며,
V0 mass는 \eqref{eq:rd-area-mass}의 homogeneous rule로 area에서 결정한다.
따라서 fixed token 수가 variable Gaussian 수로 되돌아가는 readout은 명시적이지만, Gaussian 사이 physical
edge나 persistent graph를 만들지는 않는다.

\subsection{내부 relation과 teacher supervision}

Token relation \(A_{kl}\)는 relative geometry, attachment context와 learned token feature로 계산되고
같은 forward pass 안에서만 소비된다.
Training mesh의 same-surface, geodesic 및 layer label은 attention bias, contrastive retrieval 또는
teacher relation probe의 auxiliary loss로 사용할 수 있다. 그러나 \(A_{kl}\)를 physical edge로 threshold해
solver에 전달하지 않는다.

Latent token이 의미 있는지는 다음 diagnostic으로 확인한다.
\begin{itemize}
  \item independent GS variant 사이 token coverage와 response consistency
  \item token collapse/duplicate 비율과 surface coverage
  \item same-surface neighbor retrieval 및 close-layer shortcut error
  \item attachment/free-edge query에 대한 global-context ablation
  \item internal relation을 제거한 capacity-matched direct-set baseline과 response 차이
\end{itemize}
Diagnostic이 motion gain으로 이어지지 않으면 latent-anchor 자체를 기여로 주장하지 않는다.

\section{네트워크가 출력하는 response package}

\subsection{Response의 정확한 의미}

Student setup network \(F_\theta\)는 deformed Gaussian이나 next pose를 출력하지 않는다. 먼저
target-allowed immutable runtime payload를
\[
  \mathcal D_{\mathrm{rt}}^0=\left(\operatorname{id/hash}(\mathcal G^0),\mathcal G^0,
  \vartheta_{\mathrm{metric}},\vartheta_{\mathrm{mat}},\mathcal C_{\mathrm{attach}},
  \vartheta_{\mathrm{aero}},\vartheta_{\mathrm{domain}}\right)
\]
로 둔다. Teacher mesh에서 복사한 normal, validity, area 또는 correspondence는 여기에 들어갈 수 없다.
\begin{equation}
  \begin{aligned}
  \mathcal P&=\operatorname{Pack}\!\left(\mathcal D_{\mathrm{rt}}^0,F_\theta(\mathcal I_{\mathrm{setup}})\right)\\
  &=\left(\mathcal D_{\mathrm{rt}}^0,a,m,n^0,\gamma_{\mathrm{valid}},\gamma_{\mathrm{ori}},
  B_G,\Omega_G^{\mathrm{base}},Z_G^{\mathrm{base}},
  \{\mathcal I_k,B_{L,k,\mathcal I_k},\Omega_{L,k}^{\mathrm{base}},
  Z_{L,k}^{\mathrm{base}},c_k\}_{k=1}^{K_L},B_R\right).
  \end{aligned}
  \label{eq:rd-response-package}
\end{equation}
여기서 \(a_i\ge0\)와 \(m_i\ge0\)는 Gaussian별 surface-area 및 mass quadrature,
\(B_G\)는 Global translational response field, \(\mathcal I_k\)는 R0에서 동결한 physical-valid
정책을 통과한 free Gaussian만 포함하는 ordered Local support이고
\(B_{L,k,\mathcal I_k}\)는 그 support에만 저장하는 Local field다.
\(\Omega^{\mathrm{base}}\)는 immutable positive angular-frequency diagonal matrix이고
\(Z^{\mathrm{base}}\)는 immutable positive dimensionless modal damping-ratio diagonal matrix이며,
\(n_i^0,\gamma_{\mathrm{valid},i},\gamma_{\mathrm{ori},i}\)는 target GS에서 frozen model이 예측하거나
versioned covariance rule로 계산한 rest surface payload다. \(c_k\)는 support/confidence metadata다.
\(B_R=(B_{R,G},\{B_{R,L,k}\})\)를 쓰는 branch에서는 각 angular field가 같은 이름의 translational
coordinate를 공유한다. 이를 쓰지 않는 core branch는 transported means에서 계산하는 versioned local-affine
normal update를 package에 기록한다. 두 branch 중 하나를 R0에서 freeze하고 asset별로 바꾸지 않는다.
Network inference는 setup에서 한 번 수행하고 package를 cache한다.

여기서 physical-valid Gaussian의 정확한 index set과 threshold/reject 규칙은 R0의 미결정 설계다.
허용하는 물리 의미는 hard valid set, asset-level reject 또는 둘의 사전 선언 조합뿐이다.
Soft confidence를 area/mass, traction/load, response row나 quantitative denominator의 임의 배율로 쓰지 않는다.
이 정책이 동결되기 전에는 해당 normalization과 연산을 구현하지 않는다.

Runtime의 frame response는 이 package와 current wind로 계산한 generalized acceleration/velocity increment다.
따라서 ``response를 출력한다''는 말은 \textbf{setup network가 causal force-to-motion operator를
parameterize하고, frame evaluator가 그 operator의 response increment를 계산한다}는 뜻이다.

\subsection{Area/mass measure와 density consistency}

Gaussian은 texture와 reconstruction error에 따라 비균일하게 배치되므로 opacity나 primitive count를
surface area로 간주하지 않는다. V0 homogeneous-thickness/density core는 positive area logits만
normalize하고 mass를 declared areal density에서 파생한다.
\begin{equation}
  a_i=A_{\mathrm{ref}}\frac{\operatorname{softplus}(\xi_i^a)+\epsilon_a}
  {\sum_j(\operatorname{softplus}(\xi_j^a)+\epsilon_a)},\qquad
  m_i=\frac{M_{\mathrm{ref}}}{A_{\mathrm{ref}}}a_i.
  \label{eq:rd-area-mass}
\end{equation}
\(A_{\mathrm{ref}},M_{\mathrm{ref}}>0\)는 target-side typed metric contract가 소유하고 package에
그 값과 SI 단위를 hash한다. \(M_{\mathrm{ref}}\)가 유일한 total-mass owner다.
Teacher mesh 값으로 target package를 채우거나 material preset/network head가 별도 mass scale을
곱하거나 asset별 사후 보정하지 않는다.
식 \eqref{eq:rd-area-mass}의 합 index는 R0에서 동결할 physical-valid policy를 따른다.
Hard-valid-set 경로를 선택하면 invalid row는 area, mass, load와 response field에서 함께 exact zero여야 하며,
asset-level reject 경로를 선택하면 개별 confidence를 물리량의 soft multiplier로 임의 사용하지 않는다.
Surface split/merge에서 child ownership 합이 parent ownership과 같으면 constant-cell force/work는 보존된다.
Independent reconstruction에서는 common teacher probe에서 total area, work와 response dispersion을 측정한다.

Spatially varying thickness/density가 명시적인 후속 claim과 typed input으로 필요하다는 evidence가 생길 때만
nonuniform extension을 열어
\[
  m_i=M_{\mathrm{ref}}
  \frac{\operatorname{softplus}(\xi_i^m)+\epsilon_m}
  {\sum_j(\operatorname{softplus}(\xi_j^m)+\epsilon_m)}
\]
인 별도 mass-fraction head를 사용한다. 이 branch도 합이 \(M_{\mathrm{ref}}\)가 되며 schema/domain 전체에서
사전 선언하고 homogeneous V0와 결과 후 asset별로 섞지 않는다.

\subsection{Hard attachment와 mass whitening}

Raw response field \(\widetilde B\)는 attachment row를 exact zero projection한 뒤 mass metric으로 whitening한다.
\begin{equation}
  B_c=P_{\mathrm{attach}}\widetilde B,\qquad
  B=B_c(B_c^\top M_xB_c)^{-1/2},\qquad
  M_x=\operatorname{diag}(m_iI_3).
  \label{eq:rd-attachment-whitening}
\end{equation}
그러면 attached row는 수치오차 안에서 0이고 \(B^\top M_xB=I\)다.
Gram eigenvalue가 floor 아래면 rank를 줄이거나 package를 reject한다.
Attachment를 network input에도 넣어 response shape와 pole을 조건화하되 hard exactness는 learned penalty에 맡기지 않는다.
이 SI mass-normalized convention에서
\([B]=\mathrm{kg}^{-1/2}\), \([q]=\mathrm{kg}^{1/2}\mathrm m\),
\([Q]=\mathrm{kg}^{1/2}\mathrm m\,\mathrm s^{-2}\)이고,
같은 \(q\)를 쓰는 angular field는
\([B_R]=(\mathrm{kg}^{1/2}\mathrm m)^{-1}\)다.

\subsection{Passive Global package}

Global coordinate \(q_G\)는 다음 mass-normalized damped response를 따른다.
\begin{equation}
  \ddot q_G+2Z_G^{\mathrm{eff}}\Omega_G^{\mathrm{eff}}\dot q_G+
  (\Omega_G^{\mathrm{eff}})^2q_G=Q_G,
  \qquad u_G=B_Gq_G.
  \label{eq:rd-global-response}
\end{equation}
Angular frequency와 modal damping ratio는 positive parameterization으로 만들고 \(B_G\)는
attachment projection과 whitening을 통과한다. Free response에서는 energy가 감소하는 passive LTI core가 된다.
Learned \(\Omega_S,Z_S\)는 항상 projection/whitening을 마친 \textbf{최종 coordinate}의 pole이다.
Raw field coordinate에서 예측한 pole을 basis construction 뒤 그대로 재사용하는 경로는
oracle과 main head 모두 금지한다.
V0 runtime evaluator는 대체 가능한 integrator family가 아니라 independent damped oscillator의
\textbf{single exact zero-order-hold} update 하나로 동결한다. Mode \(j\)에 대해
\begin{equation}
  \begin{aligned}
  A_{S,j}&=
  \begin{bmatrix}
    0&1\\
    -(\omega_{S,j}^{\mathrm{eff}})^2&
    -2\zeta_{S,j}^{\mathrm{eff}}\omega_{S,j}^{\mathrm{eff}}
  \end{bmatrix},
  &b&=\begin{bmatrix}0\\1\end{bmatrix},\\
  \mathcal A_{S,j}&=\begin{bmatrix}A_{S,j}&b\\0&0\end{bmatrix},
  &\exp(\Delta t\mathcal A_{S,j})
    &=\begin{bmatrix}\Phi_{S,j}&\Gamma_{S,j}\\0&1\end{bmatrix},\\
  \begin{bmatrix}q\\\dot q\end{bmatrix}_{j}^{t}
    &=\Phi_{S,j}\begin{bmatrix}q\\\dot q\end{bmatrix}_{j}^{t-1}
      +\Gamma_{S,j}Q_{S,j}^{t}.&&
  \end{aligned}
  \label{eq:rd-exact-zoh}
\end{equation}
Load는 frame 시작에서 한 번 sample해 해당 frame 동안 hold하고 mode마다 한 번만 update한다.
Package/run manifest는 \code{integrator_id=exact_zoh_augexp_v1}과
\code{force_sample_time=frame_start_v1}을 evaluator identity로 저장한다.
같은 manifest에 \(\Delta t\), precision, \(\max_j\omega_j^{\mathrm{eff}}\Delta t\)와
typed epsilon도 저장한다.
State-dependent wind가 외부 에너지를 공급하므로 forced rollout에는 위 numerical traction guard,
timestep/frequency envelope와
long-horizon stress test가 필요하다.
% [논문 작성 참고 -- 수정 이유]
% guard가 정상 결과를 결정하면 quadratic force law가 사실상 잘린 다른 모델이 된다. 따라서 core 결과에서는
% guard를 안정화 수단으로 사용하지 않고 activation을 failure/OOD 신호로 보고한다.

Full dense \(M,C,K\)나 nonlinear \(M(q),C(q),V(q)\)를 V0 output으로 두지 않는다.
이들은 latent gauge와 conditioning을 어렵게 하고 NeuROK/LaGSplat의 비선형 latent solver 범위와 겹친다.
Non-proportional damping이나 mode coupling이 반복적으로 필요한 경우에만 positive block extension을 검토한다.

\subsection{Learned Local residual package}

Local은 teacher full motion을 그대로 다시 학습하지 않는다. Frozen Global response가 설명한 component를 빼고
\begin{equation}
  r_L^*=r_{\mathrm{teacher}}-r_G^*
  \label{eq:rd-local-residual}
\end{equation}
로 정의한 complementary residual만 학습한다. \(r\)는 canonical surface probe의
time/frequency response이며 개별 eigenmode column이 아니다.
이 정의가 Global/Local double counting과 sign/permutation/degenerate-mode gauge 문제를 줄인다.

Locality는 soft loss가 아니라 package construction이다. 각 slot은 attached/invalid Gaussian을 제거한
ordered hard support \(\mathcal I_k=(i_{k,1},\ldots,i_{k,n_k})\)와 그 순서의 sparse field만 저장한다.
\(J_k\in\{0,1\}^{3N\times3n_k}\)를 support-to-full injection,
\(\widetilde C_k=J_k^\top\widetilde B_{L,k}\)를 raw support field,
\(M_k=J_k^\top M_xJ_k\), \(A_k=B_G^\top M_xJ_k\)라 두고
\begin{equation}
  \begin{aligned}
  \Pi_k&=I-M_k^{-1}A_k^\top
  (A_kM_k^{-1}A_k^\top)^\dagger A_k,\\
  \overline C_k&=\Pi_k\widetilde C_k,\qquad
  C_k=\overline C_k(\overline C_k^\top M_k\overline C_k)^{-1/2},\\
  B_{L,k}&=J_kC_k,\qquad
  Q_{L,k}=C_k^\top J_k^\top F .
  \end{aligned}
  \label{eq:rd-local-support-projection}
\end{equation}
즉 Global complement를 full field에 적용한 뒤 다시 mask하지 않고, support 내부에서
\(A_kC_k=0\)인 nullspace로 project한 뒤 같은 support mass metric으로 whitening한다.
따라서 construction상 support 밖 row와 attachment row가 0이고,
\(B_G^\top M_xB_{L,k}=0\), \(B_{L,k}^\top M_xB_{L,k}=I\)를 함께 보존한다.
Pseudoinverse/Gram floor 아래의 방향은 dev-frozen rule로 rank를 줄이거나 package를 reject한다.
이는 physical adjacency graph가 아니라 ordered sparse response-field payload다.

Slot 간 중복은 새 coupled solver로 제거하지 않고
\[
  G_{k\ell}=B_{L,k}^\top M_xB_{L,\ell},\qquad
  \rho_{k\ell}=\|G_{k\ell}\|_2,\qquad
  r_{\mathrm{eff}}=\operatorname{rank}_{\tau}
  [B_{L,1},\ldots,B_{L,K_L}]
\]
로 진단한다. Duplicate activation의 visible overshoot가 반복될 때만 dev에서 deterministic merge나
mutually exclusive group을 검토하며, V0에 global QR, coupled Local stiffness 또는 새 solver를 넣지 않는다.

Local은 처음부터 함께 학습하지 않는다.
\begin{enumerate}
  \item Global-only를 overfit하고 narrow-domain held-out response를 확인한다.
  \item Teacher residual atlas와 high-rank Global sweep으로 Local 필요성을 판정한다.
  \item Oracle/always-on Local이 실제 이득을 보인 뒤 Local package head를 학습한다.
  \item 마지막에 conditional selector와 temporal switching을 붙인다.
\end{enumerate}

\subsection{Conditional Local selection}

각 Local slot의 sparse generalized load로 dimensionless이며 slot rank에 정규화된 analytic score를 계산한다.
\(r_k\)가 slot rank이고 \(\omega_{L,k,j}^{\mathrm{eff}}\)가 effective angular frequency일 때 V0는
\begin{equation}
  \eta_k(t)=\eta_k^{\mathrm{qs}}(t)
  =\frac{1}{r_kM_{\mathrm{ref}}L_0^2}
  \sum_{j=1}^{r_k}
  \frac{Q_{L,k,j}(t)^2}
  {(\omega_{L,k,j}^{\mathrm{eff}})^4+\varepsilon_\omega},
  \qquad [\varepsilon_\omega]=\mathrm s^{-4},\qquad \lambda_g=0,
  \label{eq:rd-local-score}
\end{equation}
로 동결한다. \([Q^2/\omega^4]=\mathrm{kg\,m^2}\)이므로 \(\eta_k\)는 dimensionless
per-DOF scalar이며 vector \(\Omega\)에 대한 reduction이 구현마다 달라지지 않는다.
Analytic selector가 teacher/oracle-ranked desired set의 quality retention을 dev-frozen 기준으로
실패하면 V0의 Gate D를 실패 처리하고 always-on Local로 남긴다. Learned residual-risk term은
V0를 조용히 확장하는 fallback이 아니라 별도 method version이다. 그 version을 후속으로 열 경우에만
\(\eta_k=\eta_k^{\mathrm{qs}}+\lambda_g\chi_k\), \([\lambda_g]=[\chi_k]=1\)로 두고,
\(\chi_k\)는 current causal wind/load summary와 package가 소유하는 immutable descriptor
\(h_k^{\mathrm{sel}}\subset c_k\)만 받는다. Descriptor schema, checkpoint와 weight는 test 전에
동결하며 recurrent selector state나 teacher/test target을 받지 않는다.
Score는 deterministic tie-break를 갖는 desired ranking만 만들고, minimum dwell, fade와 아래
Active/Decay budget이 실제 admission을 결정한다.
Global은 항상 활성이고 Local만 budgeted된다. Selection overhead와 all-slot load projection이 전체 비용을 지배하면
selective speed claim은 철회한다.

\section{Wind coupling, response update와 Gaussian transport}

\subsection{Distributed wind load}

V0는 frame 시작의 geometry/velocity에서 한 번 sample하는 two-sided, normal-only,
normal-sign-invariant quasi-steady traction 하나만 사용한다. \(n_i^{t-1}\)은 unit current normal이고
\begin{equation}
  \begin{aligned}
  v_{\mathrm{rel},i}^{t}
    &=W_t^{\mathrm{eff}}(\mu_i^{t-1})-v_i^{t-1},
    &v_{n,i}^{t}&=(v_{\mathrm{rel},i}^{t})^\top n_i^{t-1},\\
  \tau_{\mathrm{raw},i}^{t}
    &=\kappa_{\mathrm{aero}}|v_{n,i}^{t}|v_{n,i}^{t}n_i^{t-1},
    &\tau_i^t&=\operatorname{clip}_{\|\cdot\|_2\le\tau_{\mathrm{guard}}}
    (\tau_{\mathrm{raw},i}^{t}),
    &F_i^t&=a_i\tau_i^t,\qquad Q_S^t=B_S^\top F^t .
  \end{aligned}
  \label{eq:rd-wind-pullback}
\end{equation}
% [논문 작성 참고 -- 수정 이유]
% kappa_aero [kg/m^3]는 0.5*rho_air*C_n을 합친 effective scale이다. 같은 공기 조건에서는 두 인자를
% 따로 추정할 수 없고, g_W까지 함께 바꾸면 F가 kappa_aero*g_W^2에 비례해 더 강하게 confound된다.
% V0는 response distillation이 기여이지 새로운 고정밀 aerodynamics가 기여가 아니므로 normal load만 남긴다.
% Gaussian Swaying은 drag/friction/lift와 C_D,C_F,C_L을 모두 쓰지만, 그 계수는 representative value로 고정한다.
% 접선항은 normal-only가 사전 동결한 Local/flutter 기준을 실패할 때만 새 law/version의 ablation으로 연다.
% 논문 citation 후보: \cite{yan2026gaussianswaying,diener2009windbasis}.
\(n\mapsto-n\)이면 \(v_n\mapsto-v_n\)이므로
\(|v_n|v_n n\)과 전체 traction은 변하지 않는다. Vector-norm guard를
area multiplication 전에 적용하고 force-level clamp와 asset별 scale tuning은 사용하지 않는다.
\(\kappa_{\mathrm{aero}}>0\)는 domain-wide frozen scalar이며 runtime artist control이 아니다.
정량 core에서 \(\tau_{\mathrm{guard}}\) activation은 0이어야 하고, activation frame을 정상 sample로 숨기지 않고
predeclared OOD/failure로 집계한다. Normal-only가 Local/flutter fidelity의 사전 동결 기준을 실패할 때만
\(v_{\parallel}=v_{\mathrm{rel}}-v_n n\)에 대해
\(r_{\parallel}\|v_{\parallel}\|_2v_{\parallel}\)를 고정 비율로 더하는 optional law를 새 dataset/method version에서
비교하며, asset별 \(r_{\parallel}\) 조정이나 별도 lift coefficient는 열지 않는다.
% [논문 작성 참고 -- 수정 이유]
% tau_guard는 논문의 물리 자유도가 아니라 catastrophic overflow를 잡는 implementation guard다.
% 접선 저항은 structural/modal damping과 식별이 섞이기 쉬우므로 기본법에서 제거하고, 필요성이 실험으로 확인될 때만
% fixed-ratio ablation으로 복원한다. 이 경계는 공력 항을 늘려 결과를 맞추는 사후 튜닝을 막는다.
\(S=G\)는 full field를, Local candidate는 \eqref{eq:rd-local-support-projection}의 sparse support
block만 사용한다. 이 runtime force는 \(S_T,S_r\) common-probe transfer를 거치지 않는다.
\(W_t^{\mathrm{eff}}=g_W(t)W_t\)이며 actual control과 load를 frame trace에 기록한다.
Opacity를 area substitute로 쓰지 않으며, floater/background와 close two-sided layer가 load를 받지 않도록
hard validity set과 confidence-triggered asset/domain reject를 둔다. Confidence를 force 배율로 쓰지 않는다.
Package/run hash는
\code{traction_law_id=two_sided_normal_quadratic_v1}과
\code{traction_guard_id=vector_norm_before_area_v1}을 method identity로 저장한다.
같은 hash에 \(\kappa_{\mathrm{aero}},\tau_{\mathrm{guard}}\), SI unit과 guard activation count도 포함한다.
% [논문 작성 참고 -- 수정 이유]
% 재현성에는 guard threshold와 activation 여부가 필요하지만, 이를 aerodynamic coefficient로 소개하면 안 된다.
% 논문 표에는 kappa_aero와 law를 보고하고 guard는 numerical-safety/validity 항목으로 분리한다.
\(g_W(t)=0\)인 zero-ambient-wind frame에서도 \(v_{\mathrm{rel},i}^t=-v_i^{t-1}\)이므로 aerodynamic drag는
남을 수 있다. 따라서 이 조건을 strict no-load/free response와 혼용하지 않고, 후자는
\(Q_S^t=0\)을 명시하고 aerodynamic evaluation을 끄는 별도 fixture로 측정한다.

Rest normal은 covariance-derived tangent cue, learned orientation confidence와 training-only teacher supervision을 사용한다.
Current normal은 angular response 또는 local affine approximation으로 갱신한다.
Gaussian covariance의 가장 작은 축이 물리 surface normal과 항상 같다고 가정하지 않으며,
orientation confidence가 낮은 target은 quantitative set에서 reject/fallback한다.

\subsection{Single response update}

Global state와 각 Local slot의
\begin{equation}
  s_{L,k}^t=(q_{L,k}^t,\dot q_{L,k}^t,\sigma_k^t,\gamma_k^t,
  d_{k,\mathrm{dwell}}^t,d_{k,\mathrm{hold}}^t),\qquad
  \sigma_k^t\in\{\mathrm{Inactive},\mathrm{Active},\mathrm{Decay}\},\quad
  0\le\gamma_k^t\le1
  \label{eq:rd-local-runtime-state}
\end{equation}
를 frame state가 명시적으로 소유한다. Active는 full \(Q_{L,k}^t\)로 exact-ZOH update하고,
Decay는 새 excitation 없이 \(Q_{L,k}^t=0\)인 homogeneous update를 계속하며,
Inactive는 \(q=\dot q=\gamma=0\)으로 update와 transport에서 제외한다.
Decay 중 재선택되면 state/phase를 zero하지 않고 Active로 복귀한다.
\(\gamma_k\)는 transport에만 적용하는 dev-frozen fade다. Load에
\(\gamma_k Q_{L,k}\)를 곱하지 않아 excitation과 transport를 이중 fade하지 않는다.
Fade, minimum dwell과 release hold는 frame count가 아니라 각각
\(T_{\mathrm{fade}},T_{\mathrm{dwell}},T_{\mathrm{hold}}\,[\mathrm s]\)의 physical duration으로
정의하며 state의 \(d_{\mathrm{dwell}},d_{\mathrm{hold}}\)는 actual \(\Delta t\)를 누적한다.
Active에서는 \(\gamma_k\)가 1로, Decay에서는 0으로 같은 frozen physical-time curve를 따라간다.
세 duration, fade curve/version과 discretization rule을 run hash에 저장한다.

\(\mathcal U_t=\{k:\sigma_k^t\in\{\mathrm{Active},\mathrm{Decay}\}\}\)라 하면 fixed budget은
Active뿐 아니라 Decay tail까지 함께 세어
\[
  |\mathcal U_t|\le K_A,\qquad \sum_{k\in\mathcal U_t}r_k\le D_A
\]
를 항상 만족해야 한다. Decay를 budget 밖에서 몰래 적분하거나 새 admission을 위해 visible tail을 즉시
버리지 않는다. Budget이 찼으면 deterministic priority/dwell rule에 따라 새 desired slot을 보류하며,
그 quality와 actual p95 비용을 그대로 Gate D에 포함한다. Decay는 \(\gamma_k=0\),
\(\dot\gamma_k=0\)이고 physical modal energy
\[
  E_k=\tfrac12\|\dot q_{L,k}\|_2^2+
  \tfrac12\|\Omega_{L,k}^{\mathrm{eff}}q_{L,k}\|_2^2
\]
가 dev-frozen SI threshold \(E_{\mathrm{release}}\,[\mathrm J]\) 아래로
\(T_{\mathrm{hold}}\) seconds 유지될 때만 Inactive로 release한다.

각 frame은 다음 순서를 정확히 한 번 수행한다.
\begin{enumerate}
  \item \(s_{t-1}\)에서 fade를 포함한 current mean, kinematic velocity, orientation과 load site를 복원한다.
  \item Frame-start bounded traction을 한 번 평가하고 Global load와 sparse Local candidate score를 계산한다.
  \item Desired ranking, dwell과 Active/Decay combined budget으로 state transition/admission을 결정한다.
  \item Global과 Active는 current load, Decay는 zero load로 같은 \(\Delta t\)의
  \eqref{eq:rd-exact-zoh}를 정확히 한 번 update한다.
  \item Transport-only \(\gamma_k\)를 한 번 update하고
  \(\dot\gamma_k^t=(\gamma_k^t-\gamma_k^{t-1})/\Delta t\)를 계산한다.
  \item 최종 response field와 정확한 kinematic velocity로 Gaussian을 transport한다.
  \item Release 조건을 판정하고 state transition, budget, score와 load를 frame trace에 기록한다.
\end{enumerate}
Fade 또는 state transition을 이유로 second force sample, hidden re-solve나 second response update를 수행하지 않는다.

Network가 매 frame absolute pose를 다시 생성하지 않기 때문에 arbitrary wind program의 causal control과
long-horizon stability를 독립적으로 검증할 수 있다.

\subsection{Renderer transport}

최소 translational transport는
\begin{equation}
  \begin{aligned}
  \mu_i^t&=\mu_i^0+B_{G,i}q_G^t+
  \sum_{k\in\mathcal U_t}\gamma_k^tB_{L,k,i}q_{L,k}^t,\\
  v_i^t&=B_{G,i}\dot q_G^t+
  \sum_{k\in\mathcal U_t}
  \left(\gamma_k^tB_{L,k,i}\dot q_{L,k}^t+
  \dot\gamma_k^tB_{L,k,i}q_{L,k}^t\right).
  \end{aligned}
  \label{eq:rd-mean-transport}
\end{equation}
\(\dot\gamma Bq\)를 빼면 rendered position의 finite-difference velocity와 wind coupling/velocity metric이
불일치하므로 생략하지 않는다. 같은 fade와 derivative convention을 angular field \(B_R\)에도 적용한다.
Flag-like disk Gaussian은 normal rotation이 시각적으로 중요하므로 같은 generalized coordinate를 사용하는
small angular field \(B_R\)를 허용한다. Rotation은 exponential map으로 적용하고 covariance eigenvalue는 보존한다.
Scale/shear, nonlinear deformation-gradient transport와 higher-order SH rotation은 translation/rotation V0가
시각적으로 실패한 증거가 있을 때만 추가한다.

Branch identity와 base/Global mechanics는 R0--R2, learned Local의 \(B_{R,L,k}\) 또는 local-affine
current-normal path는 R5, conditional fade/derivative와 다음 frame traction replay는 R6가 소유한다.
R7은 이 upstream artifact chain을 재현하고 optional appearance evidence만 추가한다.

Attachment Gaussian의 displacement/angular row는 exact zero다. Identity, proper-rigid, covariance SPD,
finite output와 visual pop fixture를 통과해야 한다. Transport 실패 frame을 quality metric에서 제거하지 않고
method failure denominator에 유지한다.

Gate B와 Gate C가 통과하고 위 mean/covariance base renderer fixture가 성공한 후에만,
R7의 optional renderer-side appearance branch를 검토한다. Analytic frame transport를 prior로 두고
\[
  \Delta R_i^t=R_i^t(R_i^0)^\top\in\mathrm{SO}(3)
\]
로 rest 대비 proper rotation을 정의한 뒤
\begin{equation}
  c_i^t=R_{\mathrm{SH}}(\Delta R_i^t)c_i^0+\Delta c_i^t,
  \qquad
  \Delta c_i^t=
  \operatorname{clip}_{\delta_c}\!\left[
  P_{\le L_{\mathrm{app}}}U_{\mathrm{app}}
  \bigl(r_\psi(\zeta_i^t)-r_\psi(\zeta_i^0)\bigr)\right]
  \label{eq:rd-optional-appearance}
\end{equation}
와 같이 bounded residual만 더한다. \(U_{\mathrm{app}}\)는 dev-frozen low-rank basis,
\(P_{\le L_{\mathrm{app}}}\)는 degree-limited projection이며 \(\delta_c\)와 \(L_{\mathrm{app}}\)를 test 전에
동결한다. \(\zeta_i^t\)는 freeze된 response coordinate/velocity, \(\Delta R_i^t\)와 immutable
GS/material cue만 포함하며 camera pose, view direction, raster target을 받지 않는다.
\(\zeta_i^0\)는 동일한 static cue에서 \(\Delta R_i=I\), \(q_S=\dot q_S=0\)으로 둔 feature이다.
동일한 zero-state feature의 차를 취하므로 residual은 construction상 정확히 0이다. 이 readout은
stateless이며 previous residual, residual history 또는 숨은 appearance state를 입력하지 않는다.

이 branch를 학습할 때 response network, response package와 deterministic evaluator를 freeze하고
appearance gradient를 dynamics로 보내지 않는다. Runtime graph, 별도 dynamic state 또는 second response update를
추가하지 않는 renderer readout이다. Appearance 개선은 Gate B/C의 motion error를 대체하거나
실패한 motion Gate를 구제할 수 없다.

\section{학습 objective와 단계적 training}

\subsection{Mode column이 아니라 response를 맞춘다}

Teacher eigenmode는 sign, order와 repeated-eigenvalue subspace rotation이 비유일하다.
따라서 mode-by-mode \(L_2\)를 primary label로 쓰지 않는다.
Random impulse, distributed wind와 frequency probe \(p\)에 대해 canonical surface에서
\begin{equation}
  \mathcal L_{\mathrm{resp}}
  =\sum_p\left(
  \lambda_x\|u_p-u_p^*\|_{W_A}^2+
  \lambda_v\|v_p-v_p^*\|_{W_M}^2+
  \lambda_f\mathcal E_{\mathrm{FRF}}(p)+
  \lambda_T\mathcal E_{\mathrm{rollout}}(p)
  \right)
  \label{eq:rd-response-loss}
\end{equation}
를 최소화한다. \(W_A,W_M\)는 frozen common-valid probe mask 위의
\(W_{A,P},W_{M,P}\)다. 합/평균, reference-scale normalization, sequence reduction과
FRF phase mask는 R2 개발 문서에서 학습 전에 동결하고 hash한다.
FRF loss는 amplitude와 phase를 함께 보고, time loss는 position보다 velocity를 primary로 둔다.

\subsection{Loss family}

\begin{longtable}{L{0.24\linewidth}L{0.62\linewidth}}
\toprule
Loss/constraint & 목적\\
\midrule
\endfirsthead
\toprule
Loss/constraint & 목적\\
\midrule
\endhead
Response rollout & Teacher displacement/velocity와 long-horizon recovery를 physical space에서 맞춘다.\\
Frequency/phase & Global sway와 Local flutter band의 amplitude, phase와 log-PSD를 맞춘다.\\
Measure ownership & Total area/mass, work와 load integral을 teacher quadrature에 맞춘다.\\
Multi-resplat consistency & 같은 source surface의 independent GS variants가 common probe에서 같은 package response를 내게 한다.\\
Latent relation auxiliary & Same-surface retrieval과 close-layer shortcut을 줄이되 relation graph를 output으로 만들지 않는다.\\
Global/Local ownership & Local이 fixed Global residual만 설명하고 중복 energy를 최소화한다.\\
Selection sparsity/continuity & Fixed budget, turnover와 activation pop을 제한한다.\\
Exact construction & Positive angular-frequency/modal damping ratio, attachment zero와 Global/Local mass construction은 penalty가 아니라 layer로 보장한다.\\
\bottomrule
\end{longtable}

Total objective는
\begin{equation}
  \mathcal L=\mathcal L_{\mathrm{resp}}
  +\lambda_{\mathrm{measure}}\mathcal L_{\mathrm{measure}}
  +\lambda_{\mathrm{pair}}\mathcal L_{\mathrm{pair}}
  +\lambda_{\mathrm{rel}}\mathcal L_{\mathrm{rel}}
  +\lambda_{\mathrm{local}}\mathcal L_{\mathrm{local}}
  +\lambda_{\mathrm{switch}}\mathcal L_{\mathrm{switch}}.
  \label{eq:rd-total-loss}
\end{equation}
정확한 weight는 canonical method equation이 아니라 development config다.
Search range, selection metric과 final value를 train/dev/test lifecycle에 맞춰 freeze하고 run hash에 넣는다.
Analytic-only V0에서 \(\mathcal L_{\mathrm{switch}}\)는 selector score나 state automaton을 학습하지 않는다.
Optional T6에서 frozen analytic selection 아래 Local/package continuity에만 사용할지,
아예 \(\lambda_{\mathrm{switch}}=0\)으로 둘지는 R6 development fixture에서 비교한 뒤
T6 실행과 Gate D test open 전에 동결한다.
어느 경우에도 learned selector를 V0에 암묵적으로 추가하지 않는다.

\subsection{Teacher period 기준 physical-time rollout curriculum}

Rollout horizon은 frame 수로 고정하지 않는다. Converged teacher의 impulse/chirp에서 canonical
surface probe velocity FRF를 teacher area/mass measure \(w_j\)로 aggregate한
\begin{equation}
  A_v(\omega)=\frac{\sum_j w_j\|H_{v,j}^*(\omega)\|_2^2}{\sum_jw_j}
  \label{eq:rd-teacher-period-spectrum}
\end{equation}
를 구한다. Area 또는 mass 중 primary weighting, predeclared Global frequency band,
spectral window, frequency resolution, prominence threshold와 zero-frequency exclusion 규칙을 dev에서 동결한다.
그 밴드에서 prominence를 통과한 가장 강한 nonzero peak를 \(\omega_{\mathrm{dom}}\)으로 선택해
\(T_{\mathrm{dom}}=2\pi/\omega_{\mathrm{dom}}\)을 정의하고,
\(H_s=\rho_sT_{\mathrm{dom}}\)인 짧은 physical horizon에서 시작해 최종 long-horizon multiple까지
\(\rho_s\)를 늘린다. \(\Delta t\)가 다르면 frame 수는 달라지되 비교하는 물리 시간과 period 수는
같게 유지한다. \(\omega_{\mathrm{dom}}\), \(\rho_s\), stage transition과 final horizon은 dev에서 동결한다.

Peak availability와 Gate 판정은 다음 truth table로 고정한다.
\begin{itemize}
  \item Teacher의 별도 spatial/time refinement에서 velocity FRF/spectrum 자체가 tolerance 안에
  수렴하지 않으면 Gate A\(_T\) 실패다.
  \item Spectrum이 수렴하고 clear stable peak가 있으면 위 period-normalized schedule을 사용한다.
  \item Spectrum은 수렴했지만 broad/flat/multi-peak라 clear stable peak가 없으면 teacher failure가 아니다.
  Dev에서 domain-wide로 미리 정한 fixed physical-time direct-long horizon
  \(H_{\mathrm{long}}\,[\mathrm s]\)을 처음부터 사용하는 schedule로 fallback하고,
  frame 수가 아니라 같은 seconds를 비교하며
  teacher-period curriculum을 contribution/claim에서 삭제한다.
  \item 어느 schedule이든 final supported physical horizon의 deterministic rollout이 발산하면
  Gate A\(_T\)/B 실패다.
\end{itemize}
Fallback 사용, horizon ID와 판정 branch는 manifest와 결과 표에 기록한다.

Network는 setup에서 한 번 package를 출력하므로 recurrent next-state predictor의 teacher forcing,
scheduled sampling 또는 opaque hidden-state injection을 사용하지 않는다. 각 rollout은 동일한 초기
physical state에서 cached package를 고정하고, deterministic evaluator가 자신의 generalized coordinate와
velocity를 전개하며 horizon만 점진적으로 늘린다.
비정지 또는 잘라낸 window의 시작은 \S\ref{sec:rd-initial-state-exploration}의 별도 초기화 경계를 따른다.
명시적 generalized-state 초기화 또는 앞선 wind 구간 재생과 초기 재구성 오차를 검증하기 전에는
같은 teacher 형상·속도 파일을 갖는 것만으로 시작 상태가 일치한다고 간주하지 않는다.

Curriculum 자체의 효과는 direct-long-horizon schedule과 fixed-frame schedule을 대상으로 model, data split,
probe set을 고정하고, total deterministic-evaluator step 수와 FLOPs를 primary compute budget으로 맞춘
equal-compute ablation으로 판정한다. 이 budget을 맞추기 위해 optimizer update 수와 batch
size/composition은 조정할 수 있으며, actual GPU wall time과 peak memory를 함께 보고한다.
Period-normalized long rollout의 velocity/phase/energy drift 이득이 사라지면 curriculum을
contribution에서 삭제하고 가장 단순한 schedule을 사용한다. Final supported horizon에서 in-envelope
divergence나 decay-order 붕괴가 반복되면 짧은 horizon의 수치로 대체하지 않고 Gate A\(_T\)/B
실패로 처리한다.

\subsection{학습 순서}

\begin{enumerate}
  \item \textbf{Stage T0:} separate teacher spatial/time convergence, common-probe mapping,
  K0/R1 oracle와 SI unit/force sign/exact-ZOH 검증
  \item \textbf{Stage T1:} area/mass measure와 patch-to-token representation pretraining
  \item \textbf{Stage T2:} single-object/single-GS Global overfit 및 long rollout
  \item \textbf{Stage T3:} multi-resplat paired Global training과 narrow-domain held-out 검증
  \item \textbf{Stage T4:} frozen Global residual로 always-on Local package 학습
  \item \textbf{Stage T5:} analytic-only Local selector와 state machine 구현, 통합 및 replay 평가;
  selector score와 state automaton은 학습하지 않음
  \item \textbf{Stage T6:} optional Local/package continuity fine-tuning. Global field/pole과
  Gate B checkpoint는 freeze하며 T6 안의 Global 갱신은 금지한다. Local/package가 바뀌면 기존
  Gate C/D evidence를 무효화하고 R5/R6을 다시 실행한다. 별도 method lineage에서 Global을 바꾸면
  R4/T3의 Gate B부터 새 checkpoint와 evidence를 만들고 이후 R5/R6을 다시 실행한다.
\end{enumerate}

Direct image/perceptual render loss는 Gate B/C와 base renderer가 성공한 뒤,
freeze된 motion/package에서 R7 optional appearance residual을 학습할 때만 검토한다.
처음부터 renderer loss를 넣으면 motion Jacobian과 physical response가 틀려도 영상만 맞는 shortcut이 생길 수 있다.

\section{왜 기존 방법을 그대로 쓰지 않는가}

\begin{longtable}{L{0.20\linewidth}L{0.30\linewidth}L{0.39\linewidth}}
\toprule
방법 & 그대로 적용할 때의 핵심 공백 & 본 방법이 검증하려는 차이\\
\midrule
\endfirsthead
\toprule
방법 & 그대로 적용할 때의 핵심 공백 & 본 방법이 검증하려는 차이\\
\midrule
\endhead
LaGSplat \cite{pottier2026lagsplat} & 대상 물체의 own dynamic video에서 per-system latent Lagrangian과 GS decoder를 식별한다. 관측되지 않은 mode와 force scale이 제한된다. & Target dynamic video/fine-tune 없이 static GS에서 multi-object amortized wind response를 예측하고 systematic teacher excitation으로 Local response를 제공한다.\\
NeuROK \cite{geng2026neurok} & Static mesh에서 plausible kinematic manifold를 예측하고 generic latent Lagrangian ODE를 푼다. Exact material-specific wind response가 학습 목표가 아니다. & Pose manifold보다 known wind/material/attachment의 force-to-motion response를 직접 증류한다.\\
Simplicits \cite{modi2024simplicits} & Representation-agnostic reduced simulation이지만 object별 neural field fitting과 elastic energy evaluation이 필요하다. & Narrow-domain amortized static-GS response compiler와 density consistency, selective Local Pareto를 검증한다.\\
FreeForm \cite{xiang2026freeform} & Particle/RKPM 기반 skinning eigenmode와 generalized eigensystem을 setup한다. & Target persistent particle physics graph 없이 student package를 prediction하고 distributed wind response를 학습한다.\\
Gaussian Swaying \cite{yan2026gaussianswaying} & Surface-based GS aerodynamics를 직접 구성하지만 learned measure-invariant response distillation과 selective Local이 없다. & Same wind task의 direct baseline으로 area/response consistency와 quality--latency 차이를 검증한다.\\
이전 explicit scaffold & Anchor edge/triplet과 structural/reduced package가 해석 가능하지만 target topology extraction과 solver contract가 복잡하다. & 내부 relation만 유지하고 외부 scaffold를 제거하되 response 정확도와 안정성을 teacher/old baseline에 대해 검증한다.\\
\bottomrule
\end{longtable}

PhysGaussian과 FastPhysGS는 Gaussian의 물리 transport/실행 baseline이고
\cite{xie2024physgaussian,ma2026fastphysgs}, DiffWind는 video-conditioned wind reconstruction과
particle/grid physics 계열의 비교축이다 \cite{lei2026diffwind}. Neural Modes의 nonlinear modal-subspace 학습은
response representation 선행축으로 비교하되, 본 방법의 density-aware distributed-load와 conditional Local claim을
그 구성요소의 최초성으로 포장하지 않는다 \cite{wang2024neuralmodes}.

입출력만 GS로 바꾼 NeuROK/LaGSplat adapter는 novelty가 아니라 failure baseline이다.
Transformer, latent token, \(B^\top F\), damped oscillator와 attention 자체도 method claim이 아니다.

\section{Gate, baseline과 metric}

\subsection{Gate A --- Teacher, measure와 latent representation}

\begin{itemize}
  \item \textbf{Gate A\(_T\), teacher prerequisite:} 별도 mesh/time refinement가 canonical velocity,
  tip trajectory와 velocity spectrum에서 수렴하고, dominant frequency가 available할 때만 해당 peak
  frequency도 수렴하며, K0/R1 oracle 및 single-object Global package가 training sequence와 긴
  free decay fixture를 통과한다. Converged broad/no-peak spectrum은 위 fixed-seconds fallback을 사용한다.
  \item \textbf{Gate A\(_M\), measure:} independent GS density/seed variant의 total area/mass 합은
  construction sanity로 확인하고, primary evidence는 local measure, analytic traction-field force/work와
  common-probe response dispersion이 dev-frozen tolerance 안에 있는지로 판정한다.
  \item \textbf{Gate A\(_L\), latent representation:} capacity-matched direct-set encoder보다 latent token/relation이
  response 또는 density robustness를 개선한다.
\end{itemize}
Gate A\(_T\) 실패는 learned-response pipeline을 막는다. Gate A\(_M\)만 실패하면 density claim을 철회하고
fixed reconstruction family로 제한할 수 있다. Gate A\(_L\)만 실패하면 latent-anchor contribution을 삭제하고
direct-set response model로 진행할 수 있다. 어느 경우에도 dataset을 무작정 늘리지 않는다.

\subsection{Gate B --- Learned Global response}

Object-disjoint narrow-domain test에서 Global package가 direct next-deformation network와
fixed mean/nearest-response baseline보다 velocity, phase, material ordering과 long rollout에서 유의하게 낫거나
동등하면서 더 stable해야 한다. Test threshold는 dev에서 freeze하고 result를 본 뒤 바꾸지 않는다.
실패하면 response-package hypothesis를 중단하거나 per-family compiler로 claim을 축소한다.

\subsection{Gate C --- Local necessity}

Strong Global rank/capacity sweep과 Global+teacher/oracle Local residual은
\textbf{actual matched-p95}에서 Gate C의 필수 primary 비교를 수행한다.
Same active-DOF 비교는 별도 run/table의 mandatory secondary diagnostic으로 수행하며 두 축을 한 결과로 섞지 않는다.
추가로 Global+procedural sine flutter와 Global+band-limited noise flutter를 두되,
learned Local과 같은 spatial support proposal을 사용하고, actual matched-p95 primary 비교와
same active-DOF secondary 비교를 각각 실행한다.
이 support는 static GS만 입력하는 proposal rule/network를 test open 전에 freeze해 얻으며,
teacher residual, test target 또는 test error로 선택하지 않는다. Procedural amplitude envelope와 phase는
현재 및 과거 \(\{W_\tau,Q_\tau\}_{\tau\le t}\)만의 causal function으로 계산하고 future wind/load를 보지 않는다.
Frequency band, amplitude normalization, envelope/filter, noise spectrum, support 규칙과 복수의 noise seed를
development에서 domain-wide로 사전 선언하고, 모든 seed 결과를 aggregate하여 보고하며 best seed를
고르지 않는다. Asset별 또는 test result별 tuning은 금지한다.
Local은 다음이 모두 참일 때만 core로 남긴다.
\begin{itemize}
  \item Traveling/local gust에서 residual이 teacher/mapping noise보다 크고 여러 object/GS variant에서 반복된다.
  \item Residual이 공간적으로 localized되고 target temporal band의 amplitude/phase/PSD에 영향을 준다.
  \item Global rank를 늘리는 것보다 oracle/always-on Local이 matched-p95에서 의미 있는 품질 이득을 낸다.
  \item Learned Local이 same-support/budget procedural flutter보다 held-out response의 spatial phase와
  target-band error를 유의미하게 개선한다.
\end{itemize}
실패하면 paper를 measure-consistent Global response로 축소한다.

\subsection{Gate D --- Conditional execution과 deployment}

Conditional Local이 always-on Local의 quality를 dev-frozen allowance 안에서 유지하면서 actual p95 latency를
개선하고, fixed high-rank Global 및 이전 explicit scaffold의 frozen Pareto envelope 밖의 점을 만들어야 한다.
Memory와 throughput은 mandatory secondary diagnostic으로 함께 보고한다.
\begin{equation}
  \begin{aligned}
  E_{\mathrm{cond}}-E_{\mathrm{always}}&\le\epsilon_{E,\mathrm{keep}},\\
  L_{\mathrm{cond}}-L_{\mathrm{always}}&\le-\delta_{L,\mathrm{always}},\\
  E_{\mathrm{cond}}&\le E_{\mathcal B}(L_{\mathrm{cond}})-\delta_E,\\
  \text{or}\quad L_{\mathrm{cond}}&\le L_{\mathcal B}(E_{\mathrm{cond}})-\delta_{L,\mathcal B}.
  \end{aligned}
  \label{eq:rd-pareto-gate}
\end{equation}
여기서 error는 dev에서 정한 primary response metric, latency는 actual p95다.
\(\mathcal B\)는 test 전에 동결한 모든 fixed high-rank Global과 explicit-scaffold baseline의 non-dominated
Pareto set이고, \(E_{\mathcal B}(L)\)와 \(L_{\mathcal B}(E)\)는 각각 그 set의 dev-frozen matched-axis envelope다.
\(\epsilon_{E,\mathrm{keep}},\delta_{L,\mathrm{always}},\delta_E,\delta_{L,\mathcal B}>0\)와
envelope interpolation/tolerance를 test open 전에 freeze한다.
Selection overhead가 이득을 상쇄하면 conditional claim을 철회하고 always-on Local 또는 Global-only로 끝낸다.

\subsection{최소 baseline}

\begin{enumerate}
  \item Converged high-resolution mesh teacher/reference
  \item 이전 explicit scaffold와 optimized full/same-resolution physical baseline
  \item Direct static-GS \(\rightarrow\) next deformation predictor
  \item Capacity-matched direct static-GS \(\rightarrow\) Global response package without latent relation supervision
  \item Global rank/capacity sweep
  \item Global + oracle/always-on Local
  \item Global + dev-frozen same-support/budget procedural sine/noise Local
  \item Global + conditional learned Local
  \item Equal-splat/opacity weight, no multi-resplat consistency, no relation auxiliary ablations
  \item 단일 GS 학습, 원본 GS random subsampling, 독립적인 여러 밀도 GS 학습의 대조:
  동일 teacher·물체 분할·network capacity·학습 budget에서 증강 효과를 비교하고,
  독립 multi-resplat 조건 안에서 measure/consistency 구성의 추가 효과를 분리한다.
  \item 공개 코드와 task adaptation이 가능한 direct GS physics/wind 방법 1--2개
\end{enumerate}
NeuROK/LaGSplat처럼 공개 implementation이 없거나 큰 adaptation이 필요한 방법은 conceptual comparison과
small adapter baseline의 범위를 명시하고, 무리한 full reproduction을 P0 blocker로 두지 않는다.

\subsection{지표와 분모}

Primary motion metric은 per-sequence mass/area-weighted velocity error, displacement/tip trajectory,
amplitude/phase와 target-band log-PSD다. Rendering은 silhouette, temporal flicker와 qualitative split-screen을 보조로 둔다.
System metric은 setup time, package memory, network/force/update/transport/render 각각의 p50/p95/p99와 throughput이다.
30fps 목표는 component 시간의 단순 합이나 평균 FPS만으로 판정하지 않고, 동기화를 포함한 실제 frame 시간과
33.3 ms 초과 비율을 함께 보고한다. GPU·출력 해상도·Gaussian 수·Global rank·활성 Local budget과
바람 입력 반영 지연을 명시한다. 성능 fixture의 구체적인 조건·허용 지연 분포는 실행 전에 동결하며 아직 미정이다.

각 source object를 independent outer cluster로 두고, 그 안의 GS package와 wind sequence를 nested paired comparison한다.
Declared OOD는 실행 전 입력 조건으로만 제외한다. In-envelope NaN, invalid covariance, package rank failure,
rollout divergence와 fallback-cap 초과는 고정 denominator의 method failure다.
Low object count에서는 descriptive interval과 all-object effect를 보고하고 accept-ready population claim을 하지 않는다.

\section{구현 난이도와 실패 원인}

\begin{longtable}{L{0.23\linewidth}L{0.31\linewidth}L{0.35\linewidth}}
\toprule
위험 & 왜 어려운가 & 조기 판정/대응\\
\midrule
\endfirsthead
\toprule
위험 & 왜 어려운가 & 조기 판정/대응\\
\midrule
\endhead
Teacher--GS correspondence & Independent reconstruction의 primitive는 material point가 아니다. & K0에서 canonical surface probe transport와 force sign부터 확인한다.\\
Density invariance & Gaussian count/opacity가 바뀌면 naive mass/force가 달라진다. & Multi-resplat pair, area/mass ownership과 equal-weight ablation을 먼저 실행한다.\\
Static identifiability & 같은 geometry도 material/attachment가 다르면 response가 다르다. & 물리 metadata를 입력하고 supported preset만 주장한다.\\
Latent collapse & Fixed token이 같은 patch를 중복하거나 teacher tessellation을 암기할 수 있다. & Coverage/duplicate probe, density pair, object-group split과 token-free baseline을 둔다.\\
Response gauge & Mode sign/order/subspace가 비유일하다. & Mode column 대신 physical probe response를 학습한다.\\
Local leakage & Local이 Global sway를 다시 학습해 selector 의미가 사라질 수 있다. & Frozen Global residual, support-constrained Global-null projection과 always-on oracle Gate를 사용한다.\\
Autoregressive drift & Small one-step error가 phase와 energy drift로 누적된다. & Passive pole parameterization, stable update, multi-step/free-decay loss와 stress test를 사용한다.\\
Renderer mismatch & Mean은 맞아도 covariance/normal이 틀리면 blur와 wrong wind가 생긴다. & Translation kill 후 angular transport와 identity/rigid/SPD fixture를 추가한다.\\
Generalization failure & Narrow data에서 network가 shape/template을 암기할 수 있다. & Source-object split, held-out aspect/silhouette, no fine-tune와 failure regime를 공개한다.\\
No speed gain & Small modal solve보다 all-GS projection/transport가 지배할 수 있다. & Component profiler 후에만 selective Local speed headline을 유지한다.\\
\bottomrule
\end{longtable}

구현 난이도의 핵심은 Transformer layer를 작성하는 것보다 teacher data contract, multi-resplat pairing,
response identifiability와 held-out evidence다. Champion domain으로 범위를 좁히면 성공 확률은 높아지지만,
per-scene fitting으로 바꾸지는 않는다.

\section{구현 순서와 claim 축소}

\subsection{Milestone}

\begin{enumerate}
  \item \textbf{R0 Contract:} typed I/O, SI/mass owner, privilege boundary, domain와 failure denominator 동결
  \item \textbf{R1 Teacher smoke:} strip/flag separate reference convergence, common-probe transfer,
  network-free oracle, traction/exact-ZOH/transport fixture
  \item \textbf{R2 Global overfit:} single object/package response overfit, passive rollout, exact attachment
  \item \textbf{R3 Measure consistency:} paired independent GS density/seed invariance와 Gate A
  \item \textbf{R4 Narrow-domain Global:} object-disjoint held-out Global response와 Gate B
  \item \textbf{R5 Local necessity:} residual atlas, high-rank Global, oracle/always-on Local과 Gate C
  \item \textbf{R6 Conditional Local:} fixed-budget selection, persistent state/fade와 Gate D
  \item \textbf{R7 Renderer/paper:} frozen base transport evidence 검증, optional appearance,
  champion/failure scene, baseline/ablation, claim과 venue freeze
\end{enumerate}

R1--R2가 실패하면 dataset 확대나 Local 구현을 하지 않는다. R3 실행은 생략하지 않으며,
Gate A\(_T\) 실패는 pipeline을 막고 Gate A\(_M\) 실패는 fixed reconstruction family로,
Gate A\(_L\) 실패는 capacity-matched direct-set encoder로 축소한다.
R4가 실패하면 target-mesh-free amortized claim을 중단한다.
R5가 실패하면 Global-only 논문으로 축소하고, R6가 실패하면 always-on Local로 남긴다.

\subsection{Claim 경로}

\begin{longtable}{L{0.13\linewidth}L{0.13\linewidth}L{0.17\linewidth}L{0.17\linewidth}L{0.27\linewidth}}
\toprule
Measure & Global & Local & Conditional & 허용 claim\\
\midrule
\endfirsthead
\toprule
Measure & Global & Local & Conditional & 허용 claim\\
\midrule
\endhead
Pass & Pass & Pass & Pass & Measure-consistent selective Global--Local wind-response distillation\\
Pass & Pass & Pass & Fail & Measure-consistent always-on Global--Local response; selector claim 철회\\
Pass & Pass & Fail & -- & Measure-consistent passive Global wind-response package\\
Fail & Pass & Pass & Pass/Fail & Fixed-reconstruction-family Global--Local response; density claim 철회\\
Fail & Pass & Fail & -- & Fixed-reconstruction-family passive Global response\\
Any & Fail & -- & -- & Pivot 중단 또는 explicit scaffold/teacher solver로 복귀\\
\bottomrule
\end{longtable}

Gate A\(_L\)가 실패하면 그 실패 판정을 그대로 기록하고 latent-anchor/relation contribution을 삭제한다.
Capacity-matched direct-set encoder를 R4의 main route로 넘겨 이후 Gate B/C를 판정하며,
direct-set 결과로 latent 이득 Gate A\(_L\)를 다시 판정하지 않는다.

이전 explicit scaffold는 버리는 코드가 아니라 teacher/oracle, interpretable baseline과 fallback research route다.
새 방향이 Gate를 통과하지 못하면 archive의 설계를 근거와 함께 복원할 수 있다.

\section{논문 전략}

\subsection{Champion scene과 quantitative evidence}

Hero scene은 soft strip step/recovery, stiff strip impulse/chirp, rectangular flag traveling gust,
동일 flag의 low/high-density independent GS를 기본으로 한다. 여기에 frozen model의 pennant 또는 real flag를
qualitative transfer로 추가할 수 있다.

소수 champion scene은 시각적 설득을 담당하고, novelty evidence는 별도의 object-disjoint narrow-domain table,
multi-resplat robustness, matched-p95 Global/Local Pareto와 failure case가 담당한다.
지원 domain 밖의 모든 상황을 해결할 필요는 없지만, 지원 domain을 결과 후에 소급 정의해서도 안 된다.

\subsection{Venue ladder}

현재 animation/simulation 중심 framing의 현실적 첫 목표는 SCA full paper를 통한
Computer Graphics Forum 게재다. Global/Local quality--latency와 broader visual system evidence가 강하면
Eurographics의 CGF 경로 또는 TVCG를 검토한다. Measure invariance와 response distillation이 일반 원리로 닫히고
comprehensive validation이 모이면 SIGGRAPH와 SIGGRAPH Asia의 dual track 또는 TOG가 상한이다.

AI 저널을 목표로 바꾸려면 champion-scene system이 아니라 broad object-disjoint operator learning,
sampling invariance, OOD와 learning-theoretic/architectural insight가 논문의 중심이 되어야 한다.
현재 scope-control 목표에는 CG route가 더 적합하다.

\section{Open decisions와 freeze policy}

다음은 구현 실험에서 선택하되 test 전에 freeze할 항목이다.
\begin{itemize}
  \item Latent token 수 \(K\), patch proposal과 global attention depth
  \item Global rank/pole 수, Local slot 수와 slot당 rank
  \item Area head architecture; nonuniform mass-fraction extension은 homogeneous V0 실패 evidence 후에만 재검토
  \item Response loss의 time/frequency weight와 rollout horizon
  \item Local ordered support proposal의 크기, overlap과 rank floor
  \item Analytic selector 실패 시 optional immutable descriptor와 learned residual-risk의 dev architecture
  \item R7 optional SH convention/degree, \(L_{\mathrm{app}}\), bounded residual capacity와 higher-order appearance policy
\end{itemize}

하지만 다음은 method identity라서 R0 전에 고정한다.
\begin{itemize}
  \item Target forward에 mesh/connectivity/correspondence가 들어가지 않는다.
  \item Network output은 absolute deformation이 아니라 cached response package다.
  \item Mesh vertex는 dense teacher sample이지 one-to-one inference token이 아니다.
  \item V0는 SI-only, angular frequency(rad/s), \(M_{\mathrm{ref}}\) sole mass owner와
  homogeneous \(m_i=(M_{\mathrm{ref}}/A_{\mathrm{ref}})a_i\)를 사용한다.
  \item Training/evaluation의 \(S_T,S_r\) common-probe map과 weighted adjoint는 mandatory지만
  target runtime traction은 GS에서 직접 계산한다.
  \item V0 traction은 domain-wide frozen \(\kappa_{\mathrm{aero}}\)와 two-sided normal-only quadratic law를
  사용한다. Numerical guard는 area 곱 전에 적용하되 core activation은 0이어야 하며,
  frame-start sample과 single exact-ZOH evaluator를 사용한다.
  % [논문 작성 참고 -- 수정 이유]
  % method identity에는 식별 가능한 effective scale 하나만 남기고, tangential/lift는 core claim 밖에 둔다.
  \item Physical-valid semantics와 \(B_R\) 또는 local-affine normal/covariance base branch는
  R0/R1 전에 method identity로 동결하고 R7에서 새로 선택하지 않는다.
  \item Global은 항상 활성이고 Local은 frozen Global residual만 소유한다.
  \item Attachment zero와 Local hard support/Global-null whitening은 deterministic construction이다.
  Positive/passive modal core도 같은 원칙을 따른다.
  \item V0 Local selector는 dimensionless analytic-only이고 fade는 transport에만 적용하며
  Active와 Decay를 같은 budget에 센다.
  \item Source-object group split과 independent multi-resplat pairing을 사용한다.
  \item Local은 Gate C 통과 전 core claim이 아니다.
\end{itemize}

명백한 수학/단위 오류를 제외하면 Gate evidence 전 새 physics backend, nonlinear latent Lagrangian,
learned aerodynamics, contact solver와 additional runtime stage를 추가하지 않는다.
새 stage가 필요하면 기존 stage 하나를 대체하고 reason/metric을 manifest에 남긴다.

\section{최종 요약}

새 Wind3DGS는 training mesh로 explicit scaffold를 target에 복원하는 방법이 아니다.
Mesh simulation은 response teacher이고, static GS는 target representation이다.
Patch 기반 encoder는 density가 다른 GS를 fixed/capped latent token으로 정리하고 내부 relation과 global context를 사용한다.
그 token에서 area와 그로부터 파생한 mass measure, always-on passive Global response와
complementary Local response dictionary를 구성한다.
Runtime에는 wind load가 package를 excite하고, 필요한 Local slot만 fixed budget에서 활성화된다.

성공 여부는 다음 네 질문으로 판정한다.
\begin{enumerate}
  \item 같은 surface의 independent GS density/seed가 달라도 total work와 response가 일관되는가?
  \item Target mesh, video, fine-tune 없이 held-out static GS의 material 및 attachment 조건별 wind response를 맞추는가?
  \item Strong Global보다 learned Local residual이 실제 flutter/phase/PSD를 복원하는가?
  \item Conditional Local이 same-p95 Pareto와 renderer-ready visual quality를 동시에 개선하는가?
\end{enumerate}

네 질문에 순서대로 답하기 전에는 universal meshless dynamics, arbitrary object generalization,
real-time과 TOG-level contribution을 주장하지 않는다.

\bibliographystyle{plain}
\begingroup
\sloppy
\raggedright
\bibliography{refs_response_distilled_global_local_wind_dynamics}
\endgroup

\end{document}
